\documentclass[11pt]{article}
\usepackage[T1]{fontenc}
\usepackage[utf8]{inputenc}
\usepackage{lmodern,microtype}
\usepackage[a4paper,margin=27mm,headheight=15pt]{geometry}
\usepackage{amsmath,amssymb,amsthm,mathtools}
\usepackage{booktabs,longtable,tabularx,array,enumitem}
\usepackage[dvipsnames]{xcolor}
\usepackage{fancyhdr}
\usepackage{listings}
\usepackage{xurl}
\usepackage[colorlinks=true,linkcolor=MidnightBlue,citecolor=MidnightBlue,urlcolor=MidnightBlue]{hyperref}
\usepackage[capitalise,noabbrev]{cleveref}
\definecolor{ink}{HTML}{153047}
\definecolor{soft}{HTML}{F2F5F7}
\pagestyle{fancy}\fancyhf{}
\fancyhead[L]{\small Certified perturbation and continuation in iPRA}
\fancyhead[R]{\small Research report}
\fancyfoot[C]{\thepage}
\setlength{\parindent}{0pt}\setlength{\parskip}{5pt}
\setlength{\emergencystretch}{3em}
\setcounter{tocdepth}{2}
\numberwithin{equation}{section}
\newtheorem{theorem}{Theorem}[section]
\newtheorem{lemma}[theorem]{Lemma}
\newtheorem{proposition}[theorem]{Proposition}
\newtheorem{corollary}[theorem]{Corollary}
\theoremstyle{definition}
\newtheorem{definition}[theorem]{Definition}
\newtheorem{example}[theorem]{Example}
\newtheorem{remark}[theorem]{Remark}
\newcommand{\N}{\mathbb N}\newcommand{\Z}{\mathbb Z}
\newcommand{\Q}{\mathbb Q}\newcommand{\R}{\mathbb R}\newcommand{\C}{\mathbb C}
\newcommand{\ipra}{\mathrm{iPRA}}
\newcommand{\PR}{\mathrm{PR}}
\newcommand{\abs}[1]{\left|#1\right|}
\newcommand{\norm}[1]{\left\lVert#1\right\rVert}
\newcommand{\ceil}[1]{\left\lceil#1\right\rceil}
\newcommand{\NN}{\mathbb N}\newcommand{\QQ}{\mathbb Q}
\DeclareMathOperator{\diag}{diag}
\DeclareMathOperator{\supp}{supp}
\DeclareMathOperator{\dist}{dist}
\DeclareMathOperator{\diam}{diam}
\DeclareMathOperator{\sgn}{sgn}
\lstset{basicstyle=\ttfamily\footnotesize,breaklines=true,columns=fullflexible,keepspaces=true,frame=single,rulecolor=\color{gray!30},backgroundcolor=\color{soft},showstringspaces=false}
\begin{document}
\hypersetup{pageanchor=false}
\begin{titlepage}
\vspace*{12mm}
{\color{ink}\Huge\bfseries Certified perturbation\par and continuation in iPRA\par}
\vspace{8mm}
{\Large A constructive framework for coupled computations,\par stable polynomial factors, root clusters, and linear ODEs}
\vspace{12mm}
{\large Research report prepared for Vladimir Reshetnikov\par 19 November 2026}
\vspace{12mm}
\hrule\vspace{5mm}
\textbf{Repository baseline.} ProveIt, commit\par
\nolinkurl{3ee850b4c2d650cf01263823fd5415963747c5d0}.

\textbf{Delivered status.} Full mathematical arguments for the stated finite theorems; exact-rational executable examples; Lean reference definitions, proof candidates, and an object-theory implementation plan. The new Lean files have \emph{not} been compiled in this environment. No new kernel-checked \texttt{Prf PRAAx} endpoint is claimed.

\textbf{Main advance.} Replace several remaining analytic producer assumptions by explicit rational algorithms with finite error proofs. A weighted contraction certificate handles cyclic dependencies. A polynomial compiler yields local implicit solutions and stable factor branches. A separate finite matching theorem corrects cluster transport, and a direct Euler comparison derives the local recurrence assumed by the existing ODE arithmetic backend.
\vfill
{\small This report distinguishes a theorem in first-order iPRA, a finite mathematical proof designed for that theory, and an ambient Lean statement. The distinction is part of every claim's status, rather than an implicit qualification.}
\end{titlepage}
\pagenumbering{roman}\hypersetup{pageanchor=true}
\begin{abstract}
ProveIt's current iPRA development already contains many finite-observation frameworks and a large library of actual derivations. The useful next step is therefore to supply quantitative producers and precise interfaces at identified boundaries. We develop a finite rational method for mutually dependent approximation problems. A nonnegative gain matrix and a positive weight vector certify contraction without requiring that every unweighted row sum be small. All convergence statements are replaced by inequalities between explicitly computed finite iterates.

We instantiate the method with a uniform evaluator for coded polynomials, a certified implicit-equation construction, and a convergent preconditioned factor-lifting algorithm. The latter labels whole polynomial factors through internal root collisions. We also prove a multiplicity-sensitive finite matching theorem, explain why component balance alone gives a weaker bound, and replace unstable complex lexicographic labeling by reference locators. Finally, exact Euler polygons for rational nonautonomous linear systems give the missing local recurrence, interpolation estimate, coherent-name modulus, and parameter stability needed for certified continuation. Every infinite-looking conclusion is accompanied by its finite numerical reading and its remaining Lean obligations.
\end{abstract}
\tableofcontents
\clearpage\pagenumbering{arabic}
\section{What this report adds to the current repository}\label{sec:scope}

\subsection{A source-specific starting point}

The baseline is a single pinned repository commit, not an impression formed from older packets. The main README distinguishes three checked libraries---\texttt{ConciseIPRA}, \texttt{Derivations}, and \texttt{Bridge}---from research packets whose intake did not itself verify them \cite{repo-readme}. We inspected the directory inventory, the dependency and endpoint documents, the large Derivations index, the FIBER and other claim audits, and the relevant source modules. We did not rebuild the whole repository. Statements below about existing verification report the repository's status and source signatures; the new results have the status declared on the title page.

The current collection already has forty-five foundation packets. It also has concrete internal bounded reflection, certificate tables, realizer composition, demand closure, dyadic coherence, finite nets, and numerical precision scheduling. For example, \nolinkurl{PrecisionCostDAG} actually computes rational error budgets on a numerically encoded graph, and \nolinkurl{PrecisionCostDAGSchedule} derives its uniform dyadic schedule by quantifier-free induction \cite{repo-dag}. Thus neither a new name for contract composition nor a statement that graph scheduling is wholly missing would be a useful or accurate contribution.

The following are the boundaries addressed here.

\begin{longtable}{>{\raggedright\arraybackslash}p{.23\textwidth}>{\raggedright\arraybackslash}p{.31\textwidth}>{\raggedright\arraybackslash}p{.37\textwidth}}
\toprule
Current boundary & What is already present & Contribution of this report\\
\midrule\endhead
Coupled approximation & Numeric acyclic budgets and bounded discrete demand fixed points & A numerical positive-gain certificate for cyclic error propagation, with a geometric clock and perturbation estimates. This is a different hypothesis from monotone capped demand closure.\\
Polynomial lifting & Ring identities for one Hensel step; various specialized lifting implications & A complete rational preconditioned iteration, invariant ball, convergence modulus, uniqueness, and coefficient stability for a supplied approximate coprime factorization.\\
FIBER G1 & No identified common-coarsening/multiplicity transport endpoint; ambiguous metric bound & Two precise finite interfaces: Hall matching from sound occurrence covers, and component transport with the necessary diameter factor.\\
FIBER G2 & Explicit audit finding that complex separation does not stabilize lexicographic labels & A counterexample and stable reference-locator theorem; factor labels provide a second, algebraic alternative.\\
Euler refinement & Checked Gronwall arithmetic and Cauchy summation; local geometric recurrence assumed & Direct derivation of that recurrence for rational nonautonomous linear systems, including interpolation and input perturbation.\\
\bottomrule
\end{longtable}

The sources for these statements are the FIBER claim audit, \texttt{HenselStep}, \nolinkurl{OdeRefinement}, and \nolinkurl{DiscreteGronwall} \cite{repo-fiber,repo-hensel,repo-ode,repo-gronwall}. An open source audit is not a theorem of impossibility. Conversely, correcting an interface does not automatically complete every downstream application.

\subsection{The common mathematical idea}

A difficult object is represented by finite observations: rational vectors, coefficient lists, finite covers, or polygons. An algorithm computes such an observation at a requested precision. A finite certificate verifies the domain, controls how input errors affect output errors, and supplies an explicit bound on the iteration count. The central extension here is to allow the observations to depend on one another cyclically, provided that a positive weighted gain estimate proves contraction.

This has three complementary benefits. First, the certificate can be checked by finite rational arithmetic even when the output is a non-rational object. Second, it supplies the actual producer of observations, rather than only a rule for composing producers assumed elsewhere. Third, it identifies stable outputs: an entire polynomial factor can be stable through a multiple root, while an ordered list of its individual roots need not be.

There is no claim that arbitrary mathematics is captured uniformly by this method. The input language is deliberately restricted. Polynomial codes have a uniform primitive recursive evaluator; arbitrary primitive recursive program codes do not have a total primitive recursive universal evaluator. Arbitrary analytic functions require their own fixed evaluation symbols and proved finite contracts.

\subsection{Reading the strength of a result}

We use four status distinctions throughout.
\begin{enumerate}[leftmargin=*,itemsep=2pt]
\item \textbf{Finite mathematical theorem:} proved here using finite rational identities, explicit loops, bounded search, and quantifier-free induction after Booleanization.
\item \textbf{Represented consequence:} an explanation of what those finite inequalities mean for Cauchy names, factors, or ODE solutions. No sort of all real numbers is added to iPRA.
\item \textbf{Lean reference or candidate:} a supplied definition or proof script in the ambient Lean language. It has not been compiled here.
\item \textbf{Object-theory target:} a concrete target of the form \(\Gamma\vdash[\mathrm{PRAAx}]\varphi\), with the program, invariant, and source modules identified. These new endpoints remain to be implemented and checked.
\end{enumerate}

The full mathematical proofs are the principal deliverable. The Lean plan is sufficiently explicit to make the remaining formal work local and reviewable, but it is not a substitute for running Lean's kernel. Exact-rational computation below checks examples, constants, and counterexamples; it does not establish universal theorems by testing.

\section{Arithmetic language, representations, and proof discipline}\label{sec:arithmetic}

\subsection{The object theory}

We use first-order intuitionistic arithmetic with equality and a function symbol for each primitive recursive definition. Initial symbols are zero, successor, and projections. Composition has its defining equations. A primitive recursion symbol has equations
\[
 f(0,\mathbf x)=g(\mathbf x),\qquad
 f(n+1,\mathbf x)=h(n,f(n,\mathbf x),\mathbf x).
\]
The successor axioms are \(Sx\ne0\) and \(Sx=Sy\to x=y\). Equality has reflexivity and substitution in terms and formulas. The logical rules are the usual intuitionistic introduction and elimination rules for \(\wedge,\vee,\to,\bot,\forall,\exists\), with the eigenvariable conditions on universal introduction and existential elimination. Negation abbreviates implication into falsity. There is no unrestricted excluded-middle rule.

For each quantifier-free formula \(A(n,\mathbf x)\), induction gives
\begin{equation}\label{eq:qfind}
 A(0,\mathbf x)\wedge
 \forall n\bigl(A(n,\mathbf x)\to A(n+1,\mathbf x)\bigr)
 \ \longrightarrow\ \forall n\,A(n,\mathbf x).
\end{equation}
Parameters are permitted. The formulas used below are equations or inequalities between primitive recursive terms, or finite Boolean encodings of bounded formulas. This is the ordinary iPRA presentation described by the repository; \texttt{ConciseIPRA}'s \texttt{Prf PRAAx} is its concrete formal target \cite{repo-readme,repo-concise}.

\subsection{One explicit numerical coding}

For clarity, the mathematics can use the following elementary coding. The implementation should reuse the checked repository encodings instead of proving a second library.

Pair two naturals by the Cantor polynomial
\[
 \langle a,b\rangle=\frac{(a+b)(a+b+1)}2+b.
\]
Its coordinates are recovered by bounded searches through the code inclusively, with a default zero on malformed inputs. Consecutive triangular-number intervals prove uniqueness. The division is exact in this formula. A list is either zero or
\(\mathrm{cons}(a,t)=1+\langle a,t\rangle\). A nonempty list's tail is smaller than its code. Iterating the tail function at most the original code computes length and all entries; this is bounded recursion, not an unproved general recursive traversal. Out-of-range entries are defined to be zero. Length checks occur in every certificate.

A signed integer is represented by two naturals with value \(u-v\). A rational is a triple \((u,v,d)\) with value \((u-v)/(d+1)\). Denominators are positive by construction. Equality and order are decided by cross multiplication; for instance
\[
 \frac{u-v}{d+1}\le\frac{u'-v'}{d'+1}
 \iff (u-v)(d'+1)\le(u'-v')(d+1).
\]
After expanding the signed expressions this is an inequality of natural terms. All displayed absolute values and divisions in the paper have this interpretation. Gaussian rationals are pairs of rationals. We use their rational norm
\(\abs{x+iy}_1=\abs x+\abs y\), not an implicitly computed square root.

Vectors and matrices are length-checked lists of these codes. Matrix multiplication is a bounded double sum. A polynomial is a finite list of pairs \((c,\alpha)\), where \(c\) is rational and \(\alpha\) is a length-\(d\) list of natural exponents. Its evaluation is
\[
 \mathrm{eval}(P,x)=\sum_{(c,\alpha)\in P}c\prod_{j<d}x_j^{\alpha_j}.
\]
Repeated monomials may be combined or simply retained. Degrees, lengths, dimension, coefficients, and precision are numerical inputs. This evaluator is primitive recursive because every product and sum has a numerical bound and exponentiation is primitive recursive.

\subsection{Bounded formulas and the actual induction invariant}

\begin{lemma}[Finite Booleanization]\label{lem:boolean}
For each fixed bounded formula made from rational comparisons and bounded quantifiers, there is a primitive recursive bit with value one exactly when the formula holds. Its elementary reflection laws are provable in iPRA.
\end{lemma}
\begin{proof}
An equality or cross-multiplied inequality is decidable by bounded arithmetic. Conjunction is multiplication of bits, disjunction is the complement of the product of their complements, and negation is complement. A bounded universal quantifier is a product of its bits. A bounded existential quantifier is the complement of the product of the complements. First prove product equations and bounded first-failure search invariants by QF induction. With a fixed horizon and the hypothesis that every indexed bit is one retained in the context, the induction formula is just the guarded equation saying the current prefix product equals one; the universal hypothesis is instantiated at the step index. Conversely, a product equal to one forces each indexed bit to be one, and a failed product has a computed failing index. The dual least-success scan returns an accepted bit for an existential fold. These numerical lemmas, combined with intuitionistic quantifier rules, prove literal bounded-quantifier reflection. The quantified reflection equivalence itself is never the induction formula. Parameters and numerical bounds remain variables. Formula syntax is fixed externally for this schema; no internal truth predicate for arbitrary arithmetic is used.
\end{proof}

The repository already implements this idea much more carefully, including substitution, witnesses, and the distinction between a bit and an arbitrary nonzero numerical value \cite{repo-fiber}. We reuse it.

When a proof says ``induct on the number of computed coordinates'', its formal invariant is, for example,
\[
 \prod_{i<d}\bigl[\abs{x_i^{(n)}}\le r v_i\bigr]=1,
 \qquad
 \prod_{i<d}\bigl[2^n\abs{x_i^{(n+1)}-x_i^{(n)}}\le C v_i\bigr]=1.
\]
These are quantifier-free equations between PR terms. An unbounded assertion such as ``all later approximations are close'' is never the induction formula. We first prove the estimate for arbitrary finite indices \(n,m\), then universally close those indices. Universal hypotheses are instantiated at the current indices, rather than inserted as extra quantifiers into the induction invariant.

For a rational recurrence with changing denominators, the state stores numerator and positive denominator. One may clear all denominators in each finite comparison, or store their product as another PR state. There is no need for a priori elementary bounds on bit length: iteration of a fixed rational update is already primitive recursive. A much smaller complexity class would require a separate analysis.

\subsection{What a represented value means}

\begin{definition}[A PR name, externally indexed]\label{def:name}
A rational-vector name is a fixed PR program \(a(p,x)\), an admissibility formula \(\mathrm{Adm}(x)\), and an iPRA derivation of a specified finite Cauchy law, for example
\[
 \forall p,q,x\quad \mathrm{Adm}(x)\wedge p\le q\ \longrightarrow\
 \norm{a(p,x)-a(q,x)}\le 2^{-p}.
\]
The code \(x\) is ordinary finite data; \(\mathrm{Adm}\) includes the required shape and certificate guards. The choice of program is made in the metatheory. Equality of two such names on admissible data means agreement to every requested tolerance with an explicitly given reindexing, proved through finite inequalities. Total behavior on malformed data does not assert a Cauchy contract there.
\end{definition}

For geometric convergence we use the PR rounding function
\[
 \ell(c)=\min\{j\le \max(1,\lceil c\rceil):c\le2^j\},\qquad c\in\Q_{\ge0}.
\]
The set is nonempty because \(2^j\ge j+1\), proved by quantifier-free induction. Thus \(c2^{-(p+\ell(c))}\le2^{-p}\). When \(c=0\), take \(\ell(c)=0\). This deliberately crude finite bound avoids assuming a real logarithm algorithm.

The notation ``the limit'' is only an economical interpretation of a constructed name and its finite laws. To prove a named polynomial identity, we give an arbitrarily small residual bound. To prove uniqueness, we compare two finite approximate solutions and let the requested precision be an arbitrary natural parameter. No completeness axiom for an internal sort of reals is invoked.

\subsection{Uniformity and the limits of the interface}

It is legitimate to quantify internally over polynomial codes because the evaluator just defined is PR. It is also legitimate to quantify over the entries of a fixed finite instruction set, provided its evaluation is PR. It is not legitimate to assume a total PR map \(U(e,n)\) evaluating every PR program: if it existed, \(n\mapsto U(n,n)+1\) would be a PR function whose own index yields a contradiction. This is a finite diagonal argument, not a claim that program execution cannot be represented with an explicit fuel bound. The repository has a corresponding obstruction module \cite{repo-evaluator}.

Likewise, a proof compiler must return an actual derivation from the hypotheses used by its source theorem. A theorem about the standard natural numbers, even if checked by Lean, does not alone give \(\Gamma\vdash[\mathrm{PRAAx}]\varphi\). Our formalization plan therefore separates a rational mathematical reference layer from the numerical programs, their QF invariants, and the ultimate object-language endpoints.

\section{Cyclic quantitative contracts from finite weighted certificates}
\label{sec:feedback}

The existing numerical precision-cost construction already handles a finite
acyclic dependency graph: it constructs the rational error budgets and a safe
dyadic schedule internally, with the numerical node count as a variable
\cite{repo-dag}. Replacing its graph by a cycle is not justified by the
contract-cut rule. A cycle must instead be discharged by a finite numerical
invariant. We give such a discharge, useful for coupled implicit equations,
preconditioned polynomial systems, and mutually dependent approximation
blocks. No fixed point, completed real space, or limit is assumed in the finite
theorems below.

\subsection{Data and the certificate}

Fix a numerical dimension $m$. A vector is a finite rational list of length
$m$; a matrix is a finite rational list of $m$ rows of that length. All rational
numbers have positive denominators. Let $A=(a_{ij})$ have nonnegative entries
and let $w=(w_i)$ have positive rational entries. The certificate is the finite
set of inequalities
\begin{equation}\label{eq:half-certificate}
 a_{ij}\geq0,\qquad w_i>0,\qquad
 2\sum_{j<m} a_{ij}w_j\leq w_i\quad(i<m).
\end{equation}
Rational weights can be multiplied by a common positive denominator, so
positive integer weights lose no generality. The checker therefore needs only
integer comparisons, rational arithmetic, and bounded sums. Neither
eigenvalues nor an extremum over a real sphere is involved.

For rational vectors define
\[
 d_w(x,y)=\max_{i<m}\frac{|x_i-y_i|}{w_i},
\]
with value $0$ when $m=0$. This maximum is a bounded rational scan. In the
object theory it is usually more convenient to replace $d_w(x,y)\leq r$ by
the Boolean fold of the inequalities $|x_i-y_i|\leq rw_i$ for $i<m$.

Let $D$ be a rational domain given by a quantifier-free decidable predicate,
and let $T$ be one fixed primitive-recursive rational-vector program, possibly
with numerical data parameters and variable dimension. Its required local
contracts are domain preservation and the componentwise gain estimate
\begin{align}
 x\in D&\Longrightarrow T(x)\in D,\label{eq:T-domain}\\
 |T_i(x)-T_i(y)|&\leq\sum_{j<m}a_{ij}|x_j-y_j|
 \quad(x,y\in D, i<m).\label{eq:T-gain}
\end{align}
The contracts are assumptions to be discharged for the chosen program.
For the polynomial applications later in the report they are finite
coefficient and box estimates; the present theorem does not infer them from
the mere presence of a matrix in a file.

\begin{lemma}[Weighted finite contraction]\label{lem:weighted-contraction}
Under \eqref{eq:half-certificate} and \eqref{eq:T-gain},
\[
 d_w(Tx,Ty)\leq\tfrac12 d_w(x,y)\qquad(x,y\in D).
\]
More directly, if $|x_j-y_j|\leq r w_j$ for every $j<m$ and $r\geq0$,
then $|T_i(x)-T_i(y)|\leq(r/2)w_i$ for every $i<m$.
\end{lemma}
\begin{proof}
Multiplication by $a_{ij}\geq0$ preserves the input inequalities. Their finite
sum gives
\[
 |T_i(x)-T_i(y)|\leq r\sum_{j<m}a_{ij}w_j\leq\tfrac r2 w_i.
\]
This is already the desired componentwise conclusion. For the norm notation,
take $r$ to be the explicitly computed finite maximum. The empty-dimensional
case has no component assertions and is immediate.
\end{proof}

Weights matter even when all entries are nonnegative. For example,
\[
 A=\begin{pmatrix}0&3\\1/16&0\end{pmatrix},\qquad
 w=\begin{pmatrix}8\\1\end{pmatrix}
\]
satisfy the certificate although the first ordinary row sum is $3$.
The feedback is cyclic; an ordinary unweighted contraction estimate would
miss this certificate.

\subsection{Finite iteration constructs a name}

Define the orbit by the primitive recursion
\[
 x_0=c,\qquad x_{n+1}=T(x_n).
\]
Suppose $c\in D$ and a nonnegative rational $B$ satisfies
$|T_i(c)-c_i|\leq Bw_i$ for all $i<m$.

\begin{theorem}[Finite tail estimate and a primitive-recursive name]
\label{thm:feedback-name}
Every orbit element lies in $D$. For all natural $n,j$ and every $i<m$,
\begin{align}
 |x_{n+1,i}-x_{n,i}|&\leq B2^{-n}w_i,\label{eq:feedback-inc}\\
 |x_{n+j,i}-x_{n,i}|&\leq
 2B2^{-n}(1-2^{-j})w_i\leq 2B2^{-n}w_i.
 \label{eq:feedback-tail}
\end{align}
Let $W\geq1$ be an integer upper bound for every $w_i$, and choose a natural
$c_0$ with $2^{c_0}\geq2BW$. Then
$X(k)=x_{k+c_0}$ is a primitive-recursive rational name satisfying
\[
 |X_i(\ell)-X_i(k)|\leq2^{-k}\quad(k\leq\ell, i<m).
\]
Moreover $T(X(k))-X(k)$ tends to zero with the explicit bound
$B2^{-(k+c_0)}w_i$.
\end{theorem}
\begin{proof}
Domain preservation proves the first assertion by induction. The first
increment is bounded by $B$; the contraction lemma halves each succeeding
increment, proving \eqref{eq:feedback-inc} by induction on $n$.
For fixed $n$, telescope the $j$ increments and use their bounds. The finite
geometric identity
\[
 \sum_{h<j}2^{-h}=2(1-2^{-j})
\]
is proved by induction on $j$ and gives \eqref{eq:feedback-tail}. Finally put
$n=k+c_0$ and $j=\ell-k$. Since $w_i\leq W$ and $2BW\leq2^{c_0}$, the
right side is at most $2^{-k}$. The residual is the next increment. Every
statement used here is about finitely many rational values; a real completion
may subsequently interpret the constructed name, but does not produce it.
\end{proof}

There is no unbounded logarithm search in the clock. If $J$ is the least
integer at least $\max(1,2BW)$, a scan through $0,\ldots,J$ finds $c_0$
with $2^{c_0}\geq J$, because $2^J\geq J$ for $J\geq1$.
The resulting least such exponent is primitive recursive. One may instead
reuse the repository's binary-length routine. The uniformity is in the
numerical coefficients, dimension, weights, and requested precision, for one
fixed evaluator template. It is not uniform evaluation of arbitrary
primitive-recursive source codes.

An invariant box can itself be supplied by a finite certificate. If
$D=\{x:|x_i-c_i|\leq Rw_i\ (i<m)\}$, the gain estimate holds on this box,
and $|T_i(c)-c_i|\leq(R/2)w_i$, then
\[
 |T_i(x)-c_i|\leq |T_i(x)-T_i(c)|+|T_i(c)-c_i|
 \leq (R/2)w_i+(R/2)w_i=Rw_i.
\]
Thus domain preservation follows from the center residual and the gain
certificate. A strict spare margin allows subsequent rounding errors.

For a general rational domain $D$, the name need not represent a point of $D$ itself: $D=(0,1)$ and $T(x)=x/2$ give a counterexample. Identifying its completion as a fixed point in the declared domain requires a completion-closed domain and a continuous extension of $T$. The closed boxes and polynomial maps used below have these properties. The finite inequalities above need neither condition.

\subsection{Residual comparison and data perturbations}

The next result is useful even when the orbit is not explicitly used.

\begin{theorem}[A posteriori finite comparison]\label{thm:feedback-residual}
Let $x,y\in D$ and suppose
$d_w(Tx,x)\leq r$ and $d_w(Ty,y)\leq s$, with $r,s\geq0$.
Then
\[
 d_w(x,y)\leq2(r+s).
\]
More generally, let $S$ be a second map, $y\in D$, and suppose
$d_w(Sy,y)\leq s$ and $d_w(Ty,Sy)\leq\delta$.
The conclusion becomes
\begin{equation}\label{eq:perturbed-residual}
 d_w(x,y)\leq2(r+s+\delta).
\end{equation}
Only $T$ needs the contraction certificate for this comparison.
\end{theorem}
\begin{proof}
For $m=0$ all discrepancies vanish. Otherwise the bounded rational maximum
$E=d_w(x,y)$ is attained at a computed index $i_0$. Three triangle
inequalities, with one extra term in the perturbed case, give
\[
 E\leq r+\tfrac12 E+s+\delta.
\]
Here division by $w_{i_0}$ is permitted because it is a positive rational
from the certificate. Subtracting $E/2$ proves
$E\leq2(r+s+\delta)$. No maximum of a represented real function and no
excluded-middle decision about real equality occurs.
\end{proof}

Applied to the names of Theorem~\ref{thm:feedback-name}, this finite estimate
gives uniqueness at the represented level. To spell out the passage, let
two rational names have residual bounds converging to zero with supplied
primitive-recursive moduli. At any requested precision, choose the two
indices where $2(r+s)$ is below that precision; the displayed finite theorem
bounds their distance. Their names are therefore equivalent. With parameter
perturbations the same argument gives a computable continuity modulus for
the solution, provided the same domain and weight certificate remain valid.
This assertion is local to the certified branch. It neither selects a branch
across a singularity nor claims global uniqueness of a polynomial system.

\subsection{Inexact iteration and genuinely simultaneous error budgets}

Suppose two finite trajectories $u_0,\ldots,u_n$ and
$v_0,\ldots,v_n$ lie in $D$ and have local defects
\[
 d_w(u_{j+1},T(u_j))\leq\varepsilon_j,
 \qquad d_w(v_{j+1},T(v_j))\leq\eta_j.
\]
Contraction and the triangle inequality yield
\[
 E_{j+1}\leq\tfrac12E_j+\varepsilon_j+\eta_j,
 \qquad E_j=d_w(u_j,v_j).
\]
Induction gives the finite convolution bound
\begin{equation}\label{eq:feedback-inexact}
 E_n\leq 2^{-n}E_0+
 \sum_{j<n}2^{-(n-1-j)}(\varepsilon_j+\eta_j).
\end{equation}
The base case is equality. For the step, multiply the preceding bound by
$1/2$ and add the new defect; this shifts every exponent and appends precisely
the last summand. This proof handles all coupled coordinates at once.

For an exact trajectory $v$ and a rounded trajectory $u$ with the same initial
point, take $\varepsilon_j\leq\varepsilon4^{-j}$ and $\eta_j=0$.
For $n\geq1$,
\[
 \sum_{j<n}2^{-(n-1-j)}\varepsilon4^{-j}
 =\varepsilon2^{-(n-1)}\sum_{j<n}2^{-j}
 \leq4\varepsilon2^{-n}.
\]
The same upper bound holds at $n=0$. Combining this with
\eqref{eq:feedback-tail} proves
\[
 d_w(u_{n+j},u_n)\leq(2B+8\varepsilon)2^{-n}.
\]
Consequently a precision schedule $p_j=2j+c_1$ for a rounded fixed rational
kernel supplies a name with an explicit clock. More precisely, for $m>0$ and $\varepsilon>0$, if each coordinate rounding error is at most $2^{-p_j}$, choose $c_1$ by bounded search so that $2^{-c_1}\le\varepsilon\min_i w_i$. Then the weighted defect is at most $\varepsilon4^{-j}$. If $\varepsilon=0$, use exact iteration; if $m=0$, the trajectory is empty. The proof requires the rounded
iterates to stay in the domain; this must be checked from a spare invariant-box
margin or enforced by a certified projection. The tail estimate alone does
not imply domain preservation.

\subsection{A finite Neumann producer}

For the nonnegative matrix $A$ in the certificate and any rational right side
$b$, define
\[
 z_0=0,\qquad z_{n+1}=Az_n+b.
\]
These are the actual finite sums
$z_n=\sum_{j<n}A^j b$, as induction verifies. The exact residual identity is
\[
 (I-A)z_n=b-A^n b.
\]
If $|b_i|\leq Bw_i$, the contraction lemma for the linear map gives
$|(A^n b)_i|\leq B2^{-n}w_i$. Hence the recurrence produces a solution name
to $(I-A)z=b$ with an explicit residual modulus and the same finite tail
bound as above. This is an effective producer for a cyclic linear block.
It uses neither a pre-existing inverse nor a convergence theorem from
classical analysis. Arbitrary finite matrix dimensions are handled by the
same bounded rational sum program.

\subsection{Object-theory induction and the Lean boundary}

All row estimates, domain tests, and coordinate comparisons have a decidable
rational meaning. Clear their positive denominators and encode a bounded
universal assertion by a primitive-recursive violation sum. For example the
increment invariant is the equation
\[
 \sum_{i<m}\bigl[\,2^n|x_{n+1,i}-x_{n,i}|>Bw_i\,\bigr]=0,
\]
where the bracket denotes the rational comparison bit after denominator
clearing. This is one quantifier-free equation, even though $m,n$ are numerical
variables. The tail proof uses the corresponding bounded fold over the
coordinate inequalities; the finite geometric identity and convolution
recurrence use ordinary quantifier-free induction. If domain preservation
is needed simultaneously, pair its Boolean invariant with the increment
invariant in the coded state.

The accompanying \nolinkurl{FeedbackReference.lean} gives definitions and
fully written candidate proofs for weighted domination, contraction, finite
increments, finite tails, bounded-maximum witnesses, and residual comparison
over Lean's rational numbers. It has not been compiled in this environment.
It is not an object-theory proof. The separate port plan specifies the
rational/vector codes and the required endpoints of the form
$\Gamma\vdash[\mathrm{PRAAx}]\varphi$. A successful compilation of the
reference file would still leave that object-theory port to be completed.
The exact-rational Python examples check cyclic linear and nonlinear
instances, including a matrix whose unweighted row sum exceeds one; those
checks support the implementation but are not proofs in iPRA.

The advance over acyclic budgets is thus specific: a finite positive-weight
certificate converts a cyclic numerical dependency into a terminating
precision-controlled producer. A cycle of unproved logical contracts remains
unproved; no global reflection or self-supporting proof rule is introduced.

\section{A certified compiler for polynomial implicit equations}\label{sec:polynomial}

The gain theorem becomes a producer when its hypotheses can be computed from the input syntax. Coded rational polynomials provide a broad, nontrivial class for which this is possible. This section gives the missing construction rather than treating a Lipschitz estimate as an unspecified oracle.

\subsection{Finite difference bounds from coefficients}

Let \(d\ge1\), and let a multi-index \(\alpha=(\alpha_0,\ldots,\alpha_{d-1})\) be a finite list of naturals. Put \(x^\alpha=\prod_{j<d}x_j^{\alpha_j}\). For a rational box with positive radii \(R_j\), define
\[
 D_{\alpha,j}(R)=
 \begin{cases}
 \alpha_j R_j^{\alpha_j-1}\displaystyle\prod_{k\ne j}R_k^{\alpha_k},&\alpha_j>0,\\
 0,&\alpha_j=0.
 \end{cases}
\]
Every exponent in an executed branch is nonnegative.

\begin{lemma}[Monomial difference]\label{lem:monomial}
If \(\abs{x_j},\abs{y_j}\le R_j\), then
\[
 \abs{x^\alpha-y^\alpha}
 \le\sum_{j<d}D_{\alpha,j}(R)\abs{x_j-y_j}.
\]
\end{lemma}
\begin{proof}
For each positive integer \(m\), the finite identity
\(s^m-t^m=(s-t)\sum_{i<m}s^{m-1-i}t^i\)
gives \(\abs{s^m-t^m}\le mR^{m-1}\abs{s-t}\). For the product of the \(d\) powers, replace the factors one at a time. The difference between consecutive products is the difference in the replaced factor times unchanged factors, each bounded by the appropriate power of \(R_k\). Summing the resulting \(d\) inequalities proves the claim. Both identities are proved by induction on a finite product length. No derivative or mean-value theorem is used.
\end{proof}

Suppose \(T:\Q^d\to\Q^d\) is given by finite coefficient lists
\(T_i(x)=\sum_\alpha t_{i,\alpha}x^\alpha\). Define the nonnegative rational matrix
\begin{equation}\label{eq:polygain}
 W_{ij}=\sum_\alpha\abs{t_{i,\alpha}}D_{\alpha,j}(R).
\end{equation}
Then the preceding lemma gives the componentwise finite estimate
\begin{equation}\label{eq:polylip}
 \abs{T_i(x)-T_i(y)}\le\sum_{j<d}W_{ij}\abs{x_j-y_j}
\end{equation}
throughout the box. Computing \(W\) requires only the finite polynomial code and rational arithmetic. It may overestimate the true gain because it ignores cancellation. Failure of the certificate therefore means only that this certificate did not establish contraction.

\subsection{The finite certificate and its checker}

Take a polynomial system \(F:\Q^d\to\Q^d\), a rational center \(x_0\), matrices \(B,C\), a positive rational \(r\), and positive rational weights \(v_i\). Define
\begin{equation}\label{eq:polyT}
 T(h)=h-BF(x_0+h),\qquad R_i=r v_i.
\end{equation}
The finite certificate checker verifies dimensions and positivity, the identity \(CB=I\), and the inequalities
\begin{equation}\label{eq:polycert}
 2\abs{T_i(0)}\le r v_i,\qquad
 2\sum_{j<d}W_{ij}v_j\le v_i\quad(i<d),
\end{equation}
where \(W\) is computed from the expanded polynomial \(T\) by \eqref{eq:polygain}. The expansion is generated by finite addition, multiplication, and powering of coefficient lists. Equality to \eqref{eq:polyT} follows from distributivity and the finite binomial identities. Alternatively, the checker may accept an expanded list accompanied by exact coefficient equality; an unchecked expansion is never trusted.

All these conditions are bounded rational comparisons. Their conjunction is one PR bit \(\operatorname{cert}(c)\), where \(c\) contains all the finite data. The bit is computable even when a candidate certificate fails. Matrices \(B,C\) need not be supplied by a special mathematical oracle: for example, one may try the inverse of the rational Jacobian at \(x_0\), computed by a bounded pivot procedure, and verify its inverse identity. That choice often works, but completeness of this search is not part of the theorem.

\begin{theorem}[Certified polynomial solver]\label{thm:polysolver}
Assume the checker accepts. Set \(h_0=0\) and \(h_{n+1}=T(h_n)\), and let
\[
 c_0=\max_{i<d}\frac{\abs{T_i(0)}}{v_i},\qquad
 A=\max_{i<d}\sum_{j<d}\abs{C_{ij}}v_j.
\]
For all natural \(n,m\), the following are finite rational inequalities:
\begin{align}
 \abs{(h_n)_i}&\le r v_i &&(i<d),\label{eq:polyball}\\
 \norm{h_{n+1}-h_n}_v&\le c_0 2^{-n},\label{eq:polystep}\\
 \norm{h_{n+m}-h_n}_v&\le 2c_0 2^{-n}(1-2^{-m})\le r2^{-n},\label{eq:polytail}\\
 \norm{F(x_0+h_n)}_\infty&\le A c_0 2^{-n}.\label{eq:polyresidual}
\end{align}
The iteration and these bounds are uniform in the numerical polynomial code, dimension, and precision. Their proofs reduce to QF induction over explicitly coded finite arrays.
\end{theorem}
\begin{proof}
By \eqref{eq:polylip} and \eqref{eq:polycert}, \(\norm{T(h)-T(k)}_v\le\tfrac12\norm{h-k}_v\). In particular, if \(\norm h_v\le r\), then
\[
 \norm{T(h)}_v\le\norm{T(0)}_v+\tfrac12\norm h_v
 \le r/2+r/2=r.
\]
This proves the invariant ball by induction. The first successive difference has norm \(c_0\). Each subsequent difference is at most half the preceding one, proving \eqref{eq:polystep}. Sum the \(m\) consecutive differences beginning at \(n\). The finite geometric identity gives \eqref{eq:polytail}; after multiplying by \(2^{n+m}\), this is an ordinary rational inequality with no infinite summation.

From the defining equation,
\(BF(x_0+h_n)=h_n-h_{n+1}\). Multiplication by the checked left inverse \(C\) yields
\(F(x_0+h_n)=C(h_n-h_{n+1})\). Rowwise triangle inequalities and \eqref{eq:polystep} give \eqref{eq:polyresidual}. Positivity of every denominator is an invariant of the rational representation. Thus every induction step is a finite algebraic implication on the stored state.
\end{proof}

This proves a strong represented consequence. Put \(V=\max(1,\max_{i<d}v_i)\). Choosing
\[
 N(p)=p+\ell\bigl(\max(1,rV,A c_0)\bigr)+2
\]
gives a PR name \(x_0+h_{N(p)}\) whose later values are within \(2^{-p-2}\) in coordinate maximum norm and whose polynomial residual is at most \(2^{-p-2}\). Polynomial evaluation on a bounded box is itself continuous with the explicit bound already proved. Therefore this name is a represented solution of \(F(x)=0\). The object theorem is the universally quantified finite statement \eqref{eq:polyball}--\eqref{eq:polyresidual}, not a quantifier over a new sort of completed real numbers.

\subsection{An automatic radius for Jacobian preconditioning}

The certificate can often be produced directly, without searching over
weights or radii. Let $J$ be the matrix of the linear coefficients of the
expanded polynomial $F(x_0+h)$. Compute and verify $B=J^{-1}$, taking
$C=J$. Then the linear part of $T(h)=h-BF(x_0+h)$ is exactly zero.
After combining equal monomials, write
\[
 T_i(h)=b_i+\sum_{|\alpha|\ge2}t_{i,\alpha}h^\alpha,
 \qquad
 M=\max\left(1,\max_i\sum_{|\alpha|\ge2}
                   |t_{i,\alpha}|\,|\alpha|\right).
\]
Set $v_i=1$ and $r=1/(2M)$. Since $r\le1$, the computed gain obeys
\[
 \sum_j W_{ij}
 =\sum_{|\alpha|\ge2}|t_{i,\alpha}|\,|\alpha|r^{|\alpha|-1}
 \le rM=1/2.
\]
Thus the single remaining numerical test
\[
                 \max_i|b_i|\le 1/(4M)
\]
produces an accepted certificate. All constants are coefficient sums and
one verified rational matrix inverse. This is a quantitative local implicit
equation theorem expressed entirely by finite arithmetic.

For example, with $F(x)=x^2-2$ and $x_0=3/2$, one has $J=3$, $B=1/3$ and
$T(h)=-1/12-h^2/3$. Here $M=1$ and $r=1/2$ pass the test. The explicit
rational iteration therefore constructs the unique positive square root
within $[1,2]$, with the finite tail and residual bounds already proved.
This example requires no bisection oracle or decision of the sign of a
represented real. In a general input, a singular rational $J$ or a failed
residual inequality simply leaves this sufficient method inapplicable.

\subsection{Uniqueness and an a posteriori estimate}

\begin{proposition}[Residual controls distance]\label{prop:polyapost}
For every rational \(z\) with \(\norm z_v\le r\), put \(\eta=\norm{z-T(z)}_v\). Then
\[
 \norm{z-h_n}_v\le r2^{-n}+2\eta\qquad(n\in\N).
\]
In the represented interpretation, there is at most one zero of \(F(x_0+z)\) in the box, and an approximate zero with fixed-point residual \(\eta\) is within \(2\eta\) of the produced name.
\end{proposition}
\begin{proof}
The triangle inequality gives
\(\norm{z-h_{n+1}}_v\le\eta+\tfrac12\norm{z-h_n}_v\).
Starting with \(\norm z_v\le r\), finite induction gives the sharper bound
\(r2^{-n}+2\eta(1-2^{-n})\). Dropping its final nonnegative correction proves the display. A represented exact zero has fixed-point residual as small as any requested rational tolerance; polynomial difference bounds transfer this fact to its rational approximants. Applying the finite estimate at arbitrarily large requested precision proves equality of the names. This supplies the full meaning of uniqueness without a decision procedure for equality of arbitrary real names.
\end{proof}

For the represented uniqueness statement, use rational approximants inside the closed box. Coordinatewise clamping supplies such approximants without increasing coordinate error. Polynomial Lipschitz bounds show that clamping a name of a point in the box preserves its limiting residual. Thus the domain hypothesis used in every finite comparison is explicit.

\subsection{Dependence on parameters and independence of a name}

Let \(T_a\) be a rational polynomial family, with a single invariant box and gain certificate valid for all rational parameters \(a\) in a specified rational parameter box. Suppose the same monomial calculation proves
\[
 \norm{T_a(h)-T_b(h)}_v\le L\norm{a-b}_w
\]
on that box. Then the exact iterations, both starting at zero, satisfy
\begin{equation}\label{eq:paramstability}
 \norm{h_n(a)-h_n(b)}_v\le 2L\norm{a-b}_w.
\end{equation}
Indeed the difference satisfies \(d_{n+1}\le d_n/2+L\norm{a-b}_w\), and its initial value is zero. This is a second finite geometric induction.

If a parameter is itself given by a fixed PR Cauchy name \(a(s)\) in that box, with \(\norm{a(s)-a(t)}_w\le2^{-s}\) for \(s\le t\), define
\[
 b(p)=h_{p+\ell(4r)}\bigl(a(p+\ell(8L))\bigr).
\]
For \(p\le q\), compare \(b(p)\) with the later iteration at the same first parameter, then change the parameter. Equations \eqref{eq:polytail} and \eqref{eq:paramstability} give
\[
 \norm{b(p)-b(q)}_v
 \le r2^{-p-\ell(4r)}+2L2^{-p-\ell(8L)}
 \le 2^{-p-1}.
\]
The shifts also work when \(L=0\). The same proof compares two equivalent input names and proves independence of their particular rational representatives. If the original input approximants can leave the parameter box, coordinatewise rational clamping is a PR nonexpansive map and supplies a convenient invariant-preserving representation.

This is the sense in which the framework is more general than a single approximation algorithm: it supplies composition, perturbation, and extensionality for the actual produced observations, including simultaneous implicit equations. It does not decide whether every polynomial system admits an accepting certificate, and it does not handle a singular Jacobian merely by renaming the problem.

\subsection{Inexact local evaluation without circular error assumptions}

In applications one may replace \(T(h_n)\) by a computed \(\widetilde T_n(h_n)\). The required local error bound must be proved independently of the desired output. Suppose the stronger invariant margin \(\norm{T(0)}_v\le r/4\) holds and every evaluation error is at most \(\delta\le r/4\). Then
\[
 \norm{\widetilde T_n(h)}_v\le r/4+r/2+r/4=r
\]
on the ball. Comparing this approximate iteration with the exact one gives error at most \(2\delta\) after any number of steps. To attain total tolerance \(\epsilon\), choose
\(r2^{-N}\le\epsilon/2\) and
\(\delta\le\min(r/4,\epsilon/4)\).
This derives a local precision demand from a checked stability margin. It does not use the final accuracy statement as its own local premise.

\begin{example}[A genuinely coupled polynomial map]
Consider
\[
 T_1(x,y)=\tfrac1{10}+y+\tfrac{x^2}{100},\qquad
 T_2(x,y)=\tfrac1{20}+\tfrac{x}{32}+\tfrac{y^2}{100}.
\]
Take \(v=(4,1)\) and \(r=1/4\), so the box is \(\abs x\le1\), \(\abs y\le1/4\). Its gain matrix is
\[
 W=\begin{pmatrix}1/50&1\\1/32&1/200\end{pmatrix},\qquad
 Wv=\begin{pmatrix}27/25\\13/100\end{pmatrix}
 \le\tfrac12\begin{pmatrix}4\\1\end{pmatrix}.
\]
The residual at zero is within \((rv)/2=(1/2,1/8)\). Therefore the theorem constructs the unique solution of \((x,y)=T(x,y)\) in the box. The first unweighted row sum is greater than one, so the ordinary unweighted contraction test would fail. Both variables are updated simultaneously; the weighted certificate justifies their mutual dependence.
\end{example}

\section{Stable polynomial factors through internal root collisions}\label{sec:factors}

\subsection{The problem and the relevant boundary}

One cannot obtain stable labels for arbitrary complex roots by sorting approximate boxes lexicographically. Even when all roots are distinct and well separated in the complex metric, small changes in real coordinates can reverse that sorting. Repeated roots cause an additional difficulty for individual-root derivatives. A useful alternative is to transport a whole monic factor. Its coefficient vector can remain well conditioned even when several roots inside that factor collide.

The existing \texttt{HenselStep} module proves the ring algebra of a lifting step, while its index explicitly leaves the convergence theorem outside that module \cite{repo-hensel}. Here we give a complete convergent finite construction under an explicit rational inverse certificate. This is an Archimedean coefficient estimate; it should not be confused with an automatic proof of every modular or Puiseux lifting theorem.

\subsection{Norms and the linear inverse certificate}

Work over \(K=\Q\), or \(K=\Q(i)\) with \(\abs{a+ib}_1=\abs a+\abs b\). For a polynomial \(p=\sum_jp_jX^j\), put
\[
 \norm p_1=\sum_j\abs{p_j}_1.
\]
Finite convolution and the triangle inequality give
\(\norm{pq}_1\le\norm p_1\norm q_1\).
For a pair of correction polynomials use
\(\norm{(u,v)}=\norm u_1+\norm v_1\).

Let \(g_0,h_0\) be monic of degrees \(r,s\ge1\), and let \(f\) be monic of degree \(r+s\). Corrections satisfy \(\deg u<r\), \(\deg v<s\). Define the square linear map
\begin{equation}\label{eq:sylvester-map}
 \mathcal L(u,v)=h_0u+g_0v.
\end{equation}
Its input and output coefficient spaces both have dimension \(r+s\). Its matrix is formed directly by shifted copies of the two coefficient lists. A certificate is a rational matrix \(S\) satisfying
\(S\mathcal L=\mathcal LS=I\). Exact multiplication checks these equations.

Define
\[
 K_0=\max_j\sum_i\abs{S_{ij}}_1,\qquad K=\max(1,K_0).
\]
Then \(\norm{SP}\le K\norm P_1\). The matrix can be produced by bounded rational Gaussian elimination or by the adjugate formula with a checked nonzero determinant. Failure to find an inverse is a failed certificate. In the usual field interpretation, invertibility of \(\mathcal L\) is equivalent to coprimeness of \(g_0,h_0\); the theorem only needs the explicitly checked inverse equations.

\subsection{The convergent finite lifting algorithm}

Put \(E=f-g_0h_0\), \(\epsilon=\norm E_1\), and choose a rational radius \(\rho>0\) such that
\begin{equation}\label{eq:factorcertificate}
 K\rho\le\tfrac14,\qquad K\epsilon\le\rho/2.
\end{equation}
A convenient automatic choice is
\(\rho=1/(4K)\), with acceptance condition
\(\epsilon\le1/(8K^2)\).
The recurrence is
\begin{equation}\label{eq:factoriteration}
 (u_0,v_0)=(0,0),\qquad
 (u_{n+1},v_{n+1})=S\bigl(E-u_nv_n\bigr).
\end{equation}
All products have degree below \(r+s\), so the finite-dimensional type is preserved.

\begin{theorem}[Certified factor transport]\label{thm:factortransport}
Under \eqref{eq:factorcertificate}, recurrence \eqref{eq:factoriteration} is a PR rational algorithm. Writing \(z_n=(u_n,v_n)\), it satisfies
\begin{align}
 \norm{z_n}&\le\rho,\label{eq:factorball}\\
 \norm{z_{n+1}-z_n}&\le K\epsilon\,2^{-n},\label{eq:factorstep}\\
 \norm{z_{n+m}-z_n}&\le2K\epsilon\,2^{-n},\label{eq:factortail}\\
 \norm{f-(g_0+u_n)(h_0+v_n)}_1
 &\le L_0K\epsilon\,2^{-n},\label{eq:factorresidual}
\end{align}
where \(L_0\) is the induced column-sum norm of the explicitly computed matrix \(\mathcal L\). The represented limits \(g,h\) are monic of degrees \(r,s\), satisfy \(f=gh\), and are the unique such pair with correction norm at most \(\rho\).
\end{theorem}
\begin{proof}
For \(z=(u,v)\) in the ball,
\[
 \norm{S(E-uv)}\le K\epsilon+K\norm u_1\norm v_1
 \le\rho/2+K\rho^2\le3\rho/4.
\]
Thus the map preserves the ball with spare room. For \(z'=(u',v')\) in the same ball, write
\(uv-u'v'=(u-u')v+u'(v-v')\). Then
\[
 \norm{uv-u'v'}_1
 \le\rho\bigl(\norm{u-u'}_1+\norm{v-v'}_1\bigr)
 =\rho\norm{z-z'}.
\]
The update therefore has gain at most \(K\rho\le1/4\), and in particular at most \(1/2\). The first displacement is at most \(K\epsilon\). Finite induction and the finite geometric sum prove \eqref{eq:factorball}--\eqref{eq:factortail}.

The inverse identity gives
\(\mathcal Lz_{n+1}=E-u_nv_n\). Hence the \emph{exact} finite residual identity is
\[
 f-(g_0+u_n)(h_0+v_n)
 =E-\mathcal Lz_n-u_nv_n
 =\mathcal L(z_{n+1}-z_n).
\]
Taking its coefficient norm proves \eqref{eq:factorresidual}. An explicit precision shift makes both the coefficient tail and the residual arbitrarily small. The leading coefficients stay one because the correction arrays omit the top coefficient. Finally, every pair of exact factors in the ball is a fixed point of the same update. Its contraction estimate, or the finite a posteriori argument of the previous section, proves uniqueness of its represented coefficient name.
\end{proof}

This proof does not start by assuming a limit factorization exists. It constructs the rational coefficient names and proves their multiplication law. It is valid with numerical degrees and coefficient lists because matrix formation, inversion verification, polynomial multiplication, and iteration are all bounded arithmetic programs.

\subsection{Parameter stability and continuation of factors}

Fix the reference pair \(g_0,h_0\), inverse \(S\), and radius \(\rho\). Let \(f,\widetilde f\) both satisfy the residual test. Their iterations have the same nonlinear term and the same preconditioner. Consequently
\[
 \norm{z_{n+1}-\widetilde z_{n+1}}
 \le K\norm{f-\widetilde f}_1+\tfrac12\norm{z_n-\widetilde z_n}.
\]
Since both initial corrections are zero, finite induction gives
\begin{equation}\label{eq:factorstability}
 \norm{z_n-\widetilde z_n}\le2K\norm{f-\widetilde f}_1
\end{equation}
for every \(n\). Passing only to the represented interpretation yields the same coefficient stability bound for the produced factors. Thus real or complex coefficient names can be handled by the reindexing construction from the preceding section, as long as their approximants remain in the certified residual neighborhood.

For continuation along a finite chain of coefficient neighborhoods, the output of one stage supplies an approximate reference factorization for the next. At each stage the algorithm must verify the new inverse and residual bounds. A finite chain of accepted stages gives a proof by finite composition. There is no assertion that an arbitrary path admits such a chain with a uniform PR length, especially if the factors become non-coprime across the chosen split.

\subsection{A worked example containing two multiple roots}

Take
\[
 g_0=X^2,\qquad h_0=(X-2)^2=X^2-4X+4,
 \qquad f_\delta=X^2(X-2)^2+\delta.
\]
In coordinates \((u_0,u_1,v_0,v_1)\), the matrices are
\[
 \mathcal L=
 \begin{pmatrix}
 4&0&0&0\\-4&4&0&0\\1&-4&1&0\\0&1&0&1
 \end{pmatrix},\qquad
 S=\begin{pmatrix}
 1/4&0&0&0\\1/4&1/4&0&0\\3/4&1&1&0\\-1/4&-1/4&0&1
 \end{pmatrix}.
\]
Their products are the identity, and \(\det\mathcal L=16\). The largest column sum of \(S\) is \(K=3/2\). Thus \(\rho=1/6\), and every \(\abs\delta\le1/18\) passes the stated residual test. The first corrections are
\[
 u_1=\tfrac\delta4(1+X),\qquad
 v_1=\tfrac\delta4(3-X).
\]
In particular, \(\delta=1/100\) gives a fully numerical nontrivial example. The coefficient correction in the limit is at most
\(2K\abs\delta=3/100\), and the joint correction error at iteration \(n\) is at most \((3/100)2^{-n}\).

At \(\delta=0\), the reference polynomial has a repeated root at zero and another repeated root at two. This causes no singularity in the factor certificate because the two factors remain coprime. Individual roots within either factor may split as \(\delta\) changes, but the two monic degree-two coefficient vectors stay within the explicit Lipschitz bound \eqref{eq:factorstability}.

There is even an elementary locator interpretation for \(\abs\delta\le1/100\). Write the limiting corrections as \(u,v\), with \(\norm u_1+\norm v_1\le3/100\). A root \(z\) of \(X^2+u\) satisfies
\(\abs z^2\le(3/100)\max(1,\abs z)\), which excludes \(\abs z\ge1/2\): for \(1/2\le\abs z\le1\) the two sides are at least \(1/4\) and at most \(3/100\), while for \(\abs z\ge1\) division gives \(\abs z\le3/100\). If \(2+w\) is a root of \((X-2)^2+v\), then
\[
 \abs w^2\le(3/100)(2+\abs w).
\]
For \(t=\abs w\ge1/2\), writing \(t=1/2+s\) shows
\(t^2-(3/100)(2+t)=7/40+(97/100)s+s^2>0\), a contradiction. Thus the two factors have all their roots in the disjoint radius-\(1/2\) disks centered at zero and two. This root interpretation uses the ordinary modulus in the ambient complex field; the coefficient construction and its proof need only rational norms. In an iPRA application, a root name and its finite residual contract give the corresponding approximate locator statement.

\subsection{What has and has not been resolved}

The construction supplies stable \emph{factor} labels through internal collisions, with a numerical modulus and a verified coefficient producer. It does not derive the existence of arbitrary disjoint root boxes, certify their multiplicities from scratch, or provide a uniformly stable ordering of individual complex roots. The next section treats finite multiplicity covers and locators on their own precise terms. Combining that interface with polynomial factor information still requires the corresponding coverage/count certificate; it is not obtained from small coefficient residual alone.

\section{Multiplicity transport and stable labels}
\label{sec:cluster-transport}

The current FIBER audit identifies two specific unresolved issues: a short
edge in an incidence graph does not give the same diameter bound for its
whole component, and a metric gap between complex roots does not stabilize
lexicographic ordering of rational boxes~\cite{repo-fiber}.  We resolve the
finite mathematical issues here.  A polynomial application must still
produce a sound complete multiplicity certificate.  No matching or
certificate-independent observation is assumed in the theorems below.

\subsection{The certificate interface}

Let $\Omega=\{0,\ldots,N-1\}$ be a finite set of \emph{occurrences}; repeated
roots have separate occurrences.  Let $x_t$, $t<N$, be points of a metric
space.  A weighted enclosure consists of centers $a_0,\ldots,a_{m-1}$,
positive natural weights $r_0,\ldots,r_{m-1}$, a radius $\epsilon_A\geq0$,
and a coverage witness
\[
 \alpha:\Omega\longrightarrow\{0,\ldots,m-1\},\qquad
 |\alpha^{-1}(i)|=r_i,\qquad d(x_t,a_{\alpha(t)})\leq\epsilon_A.
 \tag{C-A}
\]
A second enclosure $(b_j,s_j,\epsilon_B,\beta)$ has the analogous properties
for the \emph{same} occurrence multiset.  The occurrence data may be private:
only the two lists of centers and weights need be public.  For boxes, choose
a rational point in each box as center; a proved diameter bound gives a
radius bound.  One may use sharper midpoint radii when available.

These are explicit completeness and multiplicity assumptions, not consequences
of the geometry of the boxes alone.  For polynomial roots they require a
root-counting, factor, or common-refinement theorem.  In the arithmetic
implementation, finite incidence tables, count equations and distance
certificates replace the displayed metric notation.  If actual roots are
represented by names, their analytic interpretation is a separate adapter;
the finite matching program does not inspect those names.

Expand the first family into the $N$ slots
\[
       A^*=\{(i,u):i<m,\ 0\leq u<r_i\},
\]
ordered by $i$ and then $u$, and expand the second family similarly.  A slot
has its cluster's center.  Put $\delta=\epsilon_A+\epsilon_B$ and join
$(i,u)$ to $(j,v)$ when $d(a_i,b_j)\leq\delta$.  Rational $\ell^1$ or
$\ell^\infty$ distances make this a decidable finite graph, without a
square-root operation.  For a general metric interpretation, a supplied
finite adjacency relation suffices if every pair of clusters containing a
common occurrence is an edge and every edge satisfies the displayed
distance bound.

\begin{lemma}[Coverage implies Hall's inequalities]
\label{lem:coverage-hall}
For every set $S$ of slots on the first side, its neighbor set satisfies
$|S|\leq|\mathcal N(S)|$.
\end{lemma}
\begin{proof}
Let $I$ consist of those cluster indices having a slot in $S$, and let $J$
be the cluster indices adjacent to at least one index in $I$.  The triangle
inequality shows that $\alpha(t)\in I$ implies $\beta(t)\in J$:
\[
 d(a_{\alpha(t)},b_{\beta(t)})
 \leq d(a_{\alpha(t)},x_t)+d(x_t,b_{\beta(t)})\leq\delta.
\]
Counting the indicated finite sets gives
\[
 |S|\leq\sum_{i\in I}r_i
   =|\alpha^{-1}(I)|
   \leq|\beta^{-1}(J)|
   =\sum_{j\in J}s_j=|\mathcal N(S)|.
\]
The last equality holds because adjacency depends on cluster indices, not on
which copy of a center is used.  These are bounded finite sums.  In
particular, the argument does not choose a matching between occurrences.
\end{proof}

\begin{theorem}[Bounded constructive matching]
\label{thm:bounded-matching}
There is a primitive recursive algorithm which, on an $N$ by $N$ Boolean
bipartite graph, returns either a perfect matching or a set $S$ with
$|\mathcal N(S)|<|S|$.  Its correctness has a proof by quantifier-free
induction after bounded-formula reflection.  Hence the two enclosures above
have an explicitly computed permutation matching with distance at most
$\epsilon_A+\epsilon_B$ at every matched pair.
\end{theorem}
\begin{proof}
Maintain a matching covering the first $k$ left vertices, initially $k=0$.
To add left vertex $k$, search for an augmenting path.  Mark $k$ on the left.
From a marked left vertex mark all its adjacent right vertices; from a marked
right vertex already in the matching mark its matched left partner.  Record
a parent and a discovery level each time a previously unmarked vertex is
marked.  There are at most $2N$ vertices, so $2N$ rounds suffice to reach
closure; a queue implementation examines at most $N$ left vertices.

If a free right vertex is found, follow the parent records back to $k$.
Discovery levels strictly decrease, so this is a path of fewer than $2N$
edges.  Flip matched and unmatched edges along it.  Every interior vertex
keeps one matched incident edge, and the endpoints acquire one, giving a
matching of size $k+1$ using only graph edges.

If no free right vertex is reached, let $S$ and $T$ be the reached left and
right sets at closure.  Every neighbor of $S$ is in $T$, and each member of
$T$ was reached along an edge from $S$, so $T=\mathcal N(S)$.  Every vertex
of $T$ is matched.  Its partner belongs to $S\setminus\{k\}$.  Conversely,
every reached left vertex other than $k$ was reached through its matched
right partner.  Thus the old matching restricts to a bijection
$T\leftrightarrow S\setminus\{k\}$, and
\[
              |\mathcal N(S)|=|T|=|S|-1.
\]
This is a displayed Hall counterexample.  Therefore Hall's inequalities rule
out failure at each of the $N$ stages.  The final size-$N$ matching is a
permutation because both sides have exactly $N$ vertices.  When $N=0$, the
empty matching is returned immediately.

All scans have the indicated bounds.  Partial matchings, marks, parents,
levels and the stage counter are finite arrays.  Their invariants assert
only bounded index conditions, partial-inverse equations and Boolean edge
tests; each is reflected to one quantifier-free formula.  The outer stage
induction and the bounded search inductions therefore stay within iPRA.
Applying Lemma~\ref{lem:coverage-hall} eliminates the failure branch and gives
the claimed distance bound.
\end{proof}

The algorithm can always inspect vertices in numerical order, so its output
is determined by the \emph{public} adjacency table.  It does not expose a
particular coverage witness $\alpha$ or $\beta$.  This does not assert that
the same permutation is selected when a different adjacency table is used:
what is invariant is the proved matching bound and every justified symmetric
observation.  Private coverage witnesses may be sealed by existential
elimination after the matching theorem is established.

\subsection{A precise observation modulus}

For two length-$N$ lists of points define the matching distance by the finite
minimum
\[
 d_{\mathrm{match}}(a,b)
   =\min_{\sigma\in\mathfrak S_N}\max_{i<N}d(a_i,b_{\sigma(i)}).
\]
The empty-list distance is zero.  This is a finite rational computation when
the underlying distances are rational.  Theorem~\ref{thm:bounded-matching}
proves $d_{\mathrm{match}}(A^*,B^*)\leq\delta$.

Suppose $F$ is invariant under permutations, and that two ordered lists whose
coordinatewise distances are at most $\eta$ have outputs at distance at most
$\omega(\eta)$.  Reorder the second list by the constructed matching and
apply that inequality.  Permutation invariance then gives
\[
                      d(F(A^*),F(B^*))\leq\omega(\delta).
\]
Both invariance and the displayed modulus must be proved for the actual
observation program.  A phrase such as ``stable observation'' does not
replace either obligation.

\begin{proposition}[Polynomial observations]
\label{prop:cluster-polynomial-observations}
Use the submultiplicative norm $\|z\|_1=|\Re z|+|\Im z|$ on Gaussian
rationals, and suppose all centers have norm at most $R$.  If a matching has
error at most $\delta$, then, for $k\geq1$,
\[
 \left\|\sum_{i<N}a_i^k-\sum_{i<N}b_i^k\right\|_1
      \leq NkR^{k-1}\delta.
\]
For $1\leq k\leq N$, the elementary symmetric polynomials satisfy
\[
          \|e_k(a)-e_k(b)\|_1
             \leq \binom{N}{k}kR^{k-1}\delta.
\]
Thus every coefficient of the associated monic product polynomial has an
explicit finite modulus, including at collisions.
\end{proposition}
\begin{proof}
Reorder $b$ by the matching.  The identity
\[
       u^k-v^k=(u-v)\sum_{j=0}^{k-1}u^jv^{k-1-j}
\]
and submultiplicativity give $\|u^k-v^k\|_1\leq kR^{k-1}\delta$.
Summing proves the first assertion.  For each $k$-element subset, subtract
the two corresponding products by changing one factor at a time.  The
telescoping sum has $k$ terms, each bounded by $R^{k-1}\delta$.
There are $\binom{N}{k}$ such subsets, giving the second assertion.
The cases $k=0$ and $k>N$ are constant observations.  No enumeration of
continuous branches is used; every product and sum is finite.
\end{proof}

\subsection{What component balancing alone proves}

Sometimes a verifier supplies only a bipartite graph on the weighted
clusters, an upper bound $\delta$ on each edge length, and equal total red
and blue multiplicity in every connected component.  This weaker
information also gives a matching, with a different quantitative bound.

\begin{proposition}[Component transport]
\label{prop:component-transport}
In a component having $v$ vertices, there is an explicit multiplicity
matching whose every paired distance is at most $(v-1)\delta$.  For $N>0$
and positive weights, the uniform bound $(2N-1)\delta$ is valid.
\end{proposition}
\begin{proof}
List the expanded red and blue slots of the component in numerical order.
Equal total weights give equally long lists; pair equal ranks.  A shortest
edge path between the centers of any paired slots has no repeated vertex,
since deleting the part between two equal vertices shortens it.  It has at
most $v-1$ edges.  The triangle inequality gives $(v-1)\delta$.
Breadth-first search constructs such paths with at most the total number of
vertices as fuel.  Positive weights imply at most $N$ red and at most $N$
blue vertices, giving the uniform bound.  All components and rank selections
are bounded finite computations.
\end{proof}

The factor cannot simply be dropped.  On the real line take red centers
\[
                   0,\quad 1/10,\quad 2
\]
and blue centers
\[
                   1,\quad 29/10,\quad 3,
\]
all of weight one, with edges at distance at most $1$.  The graph is
connected and balanced.  The first two red centers have just the blue center
$1$ as a neighbor, so they violate Hall's inequality.  In any matching one
of these red centers must be paired with $29/10$ or $3$, at distance at
least $14/5$.  Pairing in increasing order attains $14/5$, so the exact
matching distance is $14/5>1$.  This is a counterexample to inferring a
one-edge matching bound from component balance, not to
Theorem~\ref{thm:bounded-matching}, whose complete-coverage assumptions imply
Hall's inequalities and exclude this graph.

For genuine complete covers, component balance itself also follows from the
occurrence witnesses: an occurrence's two assigned cluster vertices are
adjacent, hence in the same component.  The set of occurrences assigned to
red vertices of a component is exactly the set assigned to its blue
vertices.  Counting proves balance.  This supplies the common-coarsening
argument with its precise finite content.

\subsection{Correcting complex labels}

The audit's counterexample is decisive~\cite{repo-fiber}.  For $z^2+1$, take
arbitrarily small rational $h>0$.  Boxes centered at $(h,1)$ and $(-h,-1)$,
each of coordinate half-width $2h$, enclose $i$ and $-i$.  Boxes centered at
$(-h,1)$ and $(h,-1)$ enclose the same roots at the same accuracy.  They
remain separated by a positive metric gap as $h\to0$, but lexicographic
ordering of centers or lower corners reverses.  Thus no reduction in mesh
relative to the metric gap repairs that ordering rule.

There is a simple positive replacement which names exactly the extra data
needed.

\begin{theorem}[Labels relative to stable locators]
\label{thm:stable-locators}
Let $c_0,\ldots,c_{N-1}$ be fixed rational anchors.  Suppose each anchor
has exactly one occurrence within distance $r$, these assignments form a
complete cover, and
\[
                 d(c_i,c_j)\geq2r+\gamma
                 \quad(i\ne j),\qquad \gamma>0.
\]
Let a singleton enclosure of the same occurrences have error
$\epsilon\leq\gamma/4$.  Assign each approximate center to its nearest
anchor, using a bounded minimum scan.  This gives a bijection with the
anchor indices, independent of the singleton-enclosure certificate.
Samples carrying the same label at errors $\epsilon$ and $\epsilon'$ are
within $\epsilon+\epsilon'$ of one another.
\end{theorem}
\begin{proof}
If $x$ is the occurrence assigned to $c_j$ and $a$ approximates $x$, then
\[
                     d(a,c_j)\leq r+\epsilon.
\]
For $i\ne j$, two triangle inequalities give
\[
 d(a,c_i)\geq d(c_i,c_j)-d(c_j,x)-d(x,a)
             \geq r+\gamma-\epsilon.
\]
The second lower bound exceeds the first upper bound by at least
$\gamma-2\epsilon\geq\gamma/2>0$.  Therefore $c_j$ is the unique nearest
anchor; the numerical tie-breaking rule is never used.  The two complete
singleton covers contain exactly one occurrence for each anchor, so every
label is used once.  Two samples with label $j$ approximate that same
occurrence, giving the final inequality by the triangle inequality.
\end{proof}

The labels are relative to the fixed locator certificate.  Approximate
coordinates need not be identical across certificates; the theorem proves
that their labels refer to the same occurrence and that their difference has
the displayed bound.  Taking a dyadic error schedule yields regular names
without any global complex-root ordering.  For real roots, an ordered
interval certificate with positive order gaps gives the familiar simpler
variant.  A certified gap between real parts gives another valid ordering
guard for complex roots.

\subsection{Concrete Lean and iPRA targets}

The target has numerical dimension and graph entries inside the object
language.  Encode an $N$ by $N$ graph by bits of $e<2^{N^2}$; a subset is a
mask $s<2^N$.  Its neighbor mask has bit $j$ equal to the bounded disjunction
of $\operatorname{bit}(s,i)\wedge\operatorname{edge}(e,i,j)$ over $i<N$.
Population count is a bounded sum.  Encode a permutation by $N$ base-$(N+1)$
digits and define
\[
\begin{split}
\operatorname{Match}(N,e,p)\;:\!\!\iff\;&
  (\forall i<N\; p_i<N\ \wedge\ \operatorname{edge}(e,i,p_i)=1)\\
 &{}\wedge(\forall i,j<N\;p_i=p_j\longrightarrow i=j),\\
\operatorname{Defect}(N,e,s)\;:\!\!\iff\;&
 s<2^N\ \wedge\operatorname{count}(\operatorname{neigh}(e,s))
                  <\operatorname{count}(s).
\end{split}
\]
Every displayed quantifier is bounded.  Implement the algorithm by fixed
primitive recursion over its stage, marking and path-flipping bounds.  Its
two output projections $p(N,e)$ and $s(N,e)$ should satisfy the ordinary
closed object theorem
\[
 \forall N,e\;\bigl(e<2^{N^2}\longrightarrow
    (\operatorname{Match}(N,e,p(N,e))
       \vee\operatorname{Defect}(N,e,s(N,e)))\bigr).
\]
The decoder totalizes malformed graphs; the shape guard is retained in the
correctness theorem.  A separate object theorem converts the reflected
coverage/count formulas to the negation of every defect.  Its composition
with the algorithm gives the cluster matching endpoint, with no global
choice principle.

The Lean implementation order is: (1) actual graph/array programs and
their shape equations; (2) matching and alternating-search invariants;
(3) bounded matching-or-defect; (4) finite count monotonicity for the
coverage-to-Hall adapter; (5) actual rational distance and observation
programs; (6) the locator scan and its strict comparison proof.  Reuse
\nolinkurl{Foundations/BoundedSyntax}, \nolinkurl{Foundations/GuardedTrace},
\nolinkurl{Core/SearchParameters}, bounded sums and the rational arithmetic
modules.  The final declarations must return
\texttt{Prf PRAAx []} of these explicit formulas.  A host Lean proof of
Hall's theorem by itself would not establish that endpoint.  The supplied
reference program and finite exhaustive checks are executable evidence for
the algorithm, not a completed kernel proof.

\section{A produced analytic contract: rational linear ODE transport}
\label{sec:ode-producer}

\subsection{The precise repository boundary}

The pinned source
\nolinkurl{Derivations/Lean/IPRADerivations/Foundations/OdeRefinement.lean}
proves two arithmetic consequences of an error recurrence: a finite
Gronwall bound and geometric telescoping. Its header explicitly states that
the recurrence describing the Euler polygons is a hypothesis. The corresponding
row of \nolinkurl{Derivations/README.md} also leaves the polygons, their bounds,
interpolation, and the resulting name unconstructed. Thus another abstract
Gronwall interface would not advance this boundary. We supply the missing
producer for nonautonomous linear systems, including rational coefficient
functions and a concrete special-function application.

The result is deliberately quantitative. Every estimate below is an inequality
between finite rational arrays. Its intended object-iPRA formalization uses
signed integer numerators and positive natural denominators, clears the
positive denominators, and applies quantifier-free induction to the resulting
primitive-recursive terms. The source boundary is documented in \cite{repo-ode,repo-gronwall}. The new object derivations are a formalization
milestone; the human proofs below do not assert that those derivations have
already been checked by Lean.

\subsection{Data, norms, and the exact finite program}

Fix a finite dimension $d$. A rational coefficient kernel is a particular
primitive-recursive program evaluating a matrix $B(t)$ at a rational
$t\in[0,1]$. The word ``particular'' matters: we do not assume a total
primitive-recursive evaluator for arbitrary primitive-recursive program codes.
A finite family of such kernels, with a bounded instruction interpreter, is
another admissible implementation.

For real rational vectors use the maximum norm and for matrices the maximum
absolute row-sum norm. For complex rational vectors one may instead take
$\lVert v\rVert=\max_i(|\Re v_i|+|\Im v_i|)$ and the associated row-sum
upper bound. These norms are rational-valued and obey
\[
 \lVert v+w\rVert\le\lVert v\rVert+\lVert w\rVert,\qquad
 \lVert Av\rVert\le\lVert A\rVert\lVert v\rVert,\qquad
 \lVert AB\rVert\le\lVert A\rVert\lVert B\rVert.
\]
For example, expand $(Av)_i=\sum_jA_{ij}v_j$, apply the scalar triangle
inequality, bound each $\lVert v_j\rVert$ by the finite maximum, and take
another finite maximum. The matrix-product inequality follows by applying
the same calculation twice. These proofs are finite sum inequalities.

Assume that natural numbers $K,L$ certify, at all rational arguments in the
interval,
\begin{equation}\label{eq:ode-KL}
 \lVert B(t)\rVert\le K,\qquad
 \lVert B(t)-B(s)\rVert\le L|t-s|.
\end{equation}
Let $v$ be the initial rational vector, and let $U\ge1$ be a natural number
with $\lVert v\rVert\le U$. Put
\begin{equation}\label{eq:ode-constants}
 G=3^K,\qquad Y=GU,\qquad
 H=(L+K^2)Y,\qquad C=GH=(L+K^2)G^2U.
\end{equation}
These constants are intentionally conservative. Their explicit form is more
useful here than a sharper bound involving an unconstructed exponential.

For a positive natural $n$, define the exact Euler array by
\begin{equation}\label{eq:exact-euler}
 X_{n,0}=v,\qquad
 X_{n,j+1}=X_{n,j}+\frac1n B(j/n)X_{n,j}
 \quad(0\le j<n).
\end{equation}
For rational $t<1$, put $j=\lfloor nt\rfloor$ and
\[
 X_n(t)=X_{n,j}+(t-j/n)B(j/n)X_{n,j};
 \qquad X_n(1)=X_{n,n}.
\]
This definition uses rational arithmetic, division with remainder, and a
bounded recurrence. It is a primitive-recursive program uniformly in all
its coded finite data. The integer growth of its intermediate numerators
and denominators does not change this conclusion.

\begin{lemma}[An integer geometric bound]\label{lem:ode-growth}
For positive $n$, natural $K$, and $0\le j\le n$,
\[
 (1+K/n)^j\le3^K.
\]
\end{lemma}
\begin{proof}
First prove, by induction on $j\le n$,
\[
 (1+1/n)^j\le 1+j/n+j^2/n^2.
\]
Multiplying the right side by $1+1/n$ gives
$1+(j+1)/n+(j^2+j)/n^2+j^2/n^3$.
This is at most $1+(j+1)/n+(j+1)^2/n^2$, because
$j^2/n\le j\le j+1$. At $j=n$ the displayed upper bound is $3$.
A separate induction on $K$ proves Bernoulli's inequality
$1+K/n\le(1+1/n)^K$. Raising it to the $n$th power gives
$(1+K/n)^n\le3^K$. Finally $1+K/n\ge1$ and $j\le n$.
Every line can be multiplied by a power of the positive $n$ to become
an integer polynomial inequality. These are the same arithmetic bounds
already available in \nolinkurl{Foundations/DiscreteGronwall.lean}.
\end{proof}

\begin{lemma}[Bounds produced by the Euler program]\label{lem:ode-polybound}
At every rational $s,t\in[0,1]$,
\[
 \lVert X_n(t)\rVert\le Y,\qquad
 \lVert X_n(t)-X_n(s)\rVert\le KY|t-s|.
\]
\end{lemma}
\begin{proof}
Induction in \eqref{eq:exact-euler} gives
$\lVert X_{n,j}\rVert\le(1+K/n)^jU\le GU$.
An interpolated value is a convex combination of its two endpoint values,
so it has the same bound. On the $j$th mesh interval the slope is
$B(j/n)X_{n,j}$ and has norm at most $KY$. Split the rational interval
between $s$ and $t$ at the finitely many mesh points it contains. Apply
this bound to every piece and add their lengths. This proves the second
inequality, including intervals crossing many mesh cells.
\end{proof}

\subsection{Deriving the formerly assumed refinement recurrence}

\begin{theorem}[Complete finite refinement estimate]\label{thm:ode-produced-refinement}
The polygons computed by \eqref{eq:exact-euler} satisfy
\begin{equation}\label{eq:ode-produced-refine}
 \lVert X_n(t)-X_{2n}(t)\rVert\le\frac{C}{2n}
 \quad\text{for every rational }t\in[0,1].
\end{equation}
Consequently, for $s\ge r$,
\begin{equation}\label{eq:ode-produced-cauchy}
 \lVert X_{2^r}(t)-X_{2^s}(t)\rVert
 \le C(2^{-r}-2^{-s})\le C2^{-r}.
\end{equation}
\end{theorem}
\begin{proof}
Write $\delta=1/(2n)$, $t_j=j\delta$,
$V_j=X_n(t_j)$ and $W_j=X_{2n,j}$. Define the exact coarse residual
\[
 r_j=V_{j+1}-V_j-\delta B(t_j)V_j.
\]
If $j=2i$, then $V_j=X_{n,i}$ and the definition of interpolation gives
$r_j=0$. If $j=2i+1$, put $x=X_{n,i}$. The two coarse half-steps both have
increment $\delta B(t_{j-1})x$, while
$V_j=x+\delta B(t_{j-1})x$. Substitution, with no approximation, gives
\begin{equation}\label{eq:ode-residual-identity}
 r_j=\delta\bigl((B(t_{j-1})-B(t_j))x
       -\delta B(t_j)B(t_{j-1})x\bigr).
\end{equation}
By \eqref{eq:ode-KL} and Lemma~\ref{lem:ode-polybound},
\[
 \lVert r_j\rVert\le\delta^2(L+K^2)Y=\delta^2H.
\]
Thus $e_j=\lVert V_j-W_j\rVert$ satisfies the produced recurrence
\begin{equation}\label{eq:ode-produced-recurrence}
 e_0=0,\qquad
 e_{j+1}\le(1+K\delta)e_j+\delta^2H.
\end{equation}
Iterating this inequality and using Lemma~\ref{lem:ode-growth} with $2n$
in place of $n$ yields, for every $j\le2n$,
\[
 e_j\le G\sum_{i<j}\delta^2H
 \le G(2n)\delta^2H=\frac{C}{2n}.
\]
This is exactly the point at which the repository's
\texttt{OdeRefinement.node\_bound} can be instantiated after rational
encoding and clearing denominators. We have now supplied its recurrence
premise, rather than repackaging it.

Both polygons are affine on each interval $[t_j,t_{j+1}]$ of the fine
mesh. Their difference at an intermediate point is a convex combination
of the differences at the two endpoints. The triangle inequality therefore
extends the node estimate to every rational point, proving
\eqref{eq:ode-produced-refine}. Finally add this bound for
$n=2^r,2^{r+1},\ldots,2^{s-1}$. The finite geometric sum is
$\sum_{i=r}^{s-1}2^{-i-1}=2^{-r}-2^{-s}$, proving
\eqref{eq:ode-produced-cauchy}. Its denominator-cleared proof is the
geometric telescoping already available in the repository.
\end{proof}

\subsection{A primitive-recursive name and its exact demands}

For a natural $M\ge1$, let $\ell(M)$ be the least $j\le M$ such that
$M\le2^j$, and put $\ell(0)=0$. This is bounded search. Its search interval
is sufficient because $2^M\ge M$, proved by elementary induction. In
particular, $\ell$ is primitive recursive; it is not an unbounded
``search until the requested error is small.''

\begin{corollary}[A produced regular name]\label{cor:ode-regular-name}
For rational initial data, put $c=\ell(\max(1,C))$ and
\[
 q_p(t)=X_{2^{p+c}}(t).
\]
Then $q$ is primitive recursive, and
$\lVert q_p(t)-q_q(t)\rVert\le2^{-p}$ whenever $q\ge p$.
\end{corollary}
\begin{proof}
Apply \eqref{eq:ode-produced-cauchy} and $C\le2^c$.
\end{proof}

One can also accept a fixed primitive-recursive initial name $(v_r)$ with
\[
 \lVert v_r\rVert\le U,\qquad
 \lVert v_r-v_s\rVert\le2^{-r}+2^{-s}.
\]
We write $X_n(v;t)$ when the initial vector must be displayed.

\begin{lemma}[Initial-data and coefficient stability]\label{lem:ode-stability}
For the same coefficient kernel,
\[
 \lVert X_n(v;t)-X_n(w;t)\rVert\le G\lVert v-w\rVert.
\]
If $B,\widetilde B$ both have norm at most $K$, and
$\lVert B(t)-\widetilde B(t)\rVert\le\eta$ on the interval, then
\[
 \lVert X_n^B(v;t)-X_n^{\widetilde B}(v;t)\rVert
 \le G^2U\eta.
\]
\end{lemma}
\begin{proof}
The initial-data difference obeys the homogeneous Euler recurrence, so the
product bound in Lemma~\ref{lem:ode-growth} applies directly. For the second
claim subtract the recurrences. The two terms are a previous error multiplied
by $I+B(j/n)/n$, and the new term
$(B(j/n)-\widetilde B(j/n))X_n^{\widetilde B}(v;j/n)/n$.
Hence the norm error satisfies
\[
 e_{j+1}\le(1+K/n)e_j+\eta GU/n,\qquad e_0=0.
\]
Iteration gives $e_j\le G(\eta GU)=G^2U\eta$. The common-mesh affine
interpolation argument extends both estimates off the nodes.
\end{proof}

For named initial data and an exact rational coefficient kernel choose
\begin{equation}\label{eq:ode-demands}
 c_v=\ell(4G),\qquad c_m=\ell(4\max(1,C)),\qquad
 q_p(t)=X_{2^{p+c_m}}(v_{p+c_v};t).
\end{equation}
For $q\ge p$, compare first across meshes at the same initial vector and
then across initial vectors at the finer mesh. The estimates give
\[
 \lVert q_p(t)-q_q(t)\rVert
 \le\tfrac14 2^{-p}
      +\tfrac14(2^{-p}+2^{-q})
 \le\tfrac34 2^{-p}.
\]
The backward precision demand on the initial name is therefore the affine
map $p\mapsto p+c_v$. The mesh resource is the explicit number
$2^{p+c_m}$. This fits the existing numerical precision/cost machinery.

If coefficients themselves arrive through names, use fixed rational
approximating kernels $B_r$, with the same certified $K,L$ and a proved
uniform bound
$\lVert B_r-B_s\rVert\le2^{-r}+2^{-s}$. Allocate all three sources of
error by replacing the $4$ in the mesh and initial shifts with $8$, and
use the coefficient shift $c_B=\ell(8G^2U)$. Mesh error then contributes
at most $2^{-p}/8$, and each of the two input-name comparisons at most
$2^{-p}/4$. Their sum is below $2^{-p}$. The existence of such $B_r$,
including uniform denominator guards, must be proved from the particular
coefficient representation; the stability lemma does not invent it.

\subsection{A path-clearance producer for the Gauss equation}

We now construct the bounds \eqref{eq:ode-KL} from concrete data. This is
useful for the path-continuation theorem in
\texttt{analysis/hypergeometric\_ipra\_2026-11-16/sections/10\_continuation.tex}.
That packet supplies the intended special-function interface; our producer
provides a fully explicit finite arithmetic alternative to its local
Taylor-majorant step.

Let $z(t)=u+th$ be one complex-rational line segment, and consider a singular
point $s\in\{0,1\}$. Put $a=u-s$, $E=|h|_2^2$, and
$D=\Re(\overline a h)$, where $|x+iy|_2^2=x^2+y^2$ involves no square
root. If $E=0$, set $\mu_s=|a|_2^2$. Otherwise set
\[
 \tau=\min(1,\max(0,-D/E)),\qquad
 \mu_s=|a+\tau h|_2^2.
\]

\begin{lemma}[Exact segment clearance]\label{lem:ode-clearance}
The rational number $\mu_s$ is the minimum squared distance from the
segment to $s$. The test $\mu_s>0$ succeeds exactly when the segment misses
$s$, and supplies a lower bound valid at every rational parameter $t$.
\end{lemma}
\begin{proof}
The distance polynomial is $f(t)=|a|_2^2+2Dt+Et^2$. If $E=0$ it is constant.
If $-D/E\in[0,1]$, completing the square gives
$f(t)-f(-D/E)=E(t+D/E)^2\ge0$. If $-D/E<0$, then $D>0$ and
$f(t)-f(0)=t(2D+Et)\ge0$. If $-D/E>1$, then $D<-E$, and for $t\le1$,
$f(t)-f(1)=(t-1)(2D+E(t+1))\ge0$. The boundary cases satisfy either
adjacent formula. Thus the minimum is achieved at the displayed rational
$\tau$. It is zero precisely when $a+\tau h=0$.
\end{proof}

Assume both clearance tests pass. Write $|z|_1=|\Re z|+|\Im z|$ and set
\[
 R=\max(|u|_1,|u+h|_1),\qquad T=|h|_1,\qquad
 D_*=\frac{R(1+R)}{\mu_0\mu_1}.
\]
Convexity gives $|z(t)|_1\le R$. Since
$|1/z|_1=|z|_1/|z|_2^2$, the clearances imply
\begin{equation}\label{eq:ode-inverse-bound}
 \left|\frac1{z(1-z)}\right|_1\le D_*.
\end{equation}
For two points $z,w$ on the segment, putting $d(z)=z(1-z)$ gives
\[
 d(z)-d(w)=(z-w)(1-z-w),\qquad
 |d(z)-d(w)|_1\le(1+2R)|z-w|_1.
\]
The reciprocal difference identity therefore gives
\begin{equation}\label{eq:ode-inverse-lip}
 \left|\frac1{d(z)}-\frac1{d(w)}\right|_1
 \le D_*^2(1+2R)|z-w|_1.
\end{equation}
All tests, constants, and proofs so far are rational arithmetic.

For complex-rational hypergeometric parameters $a,b,c$, put
\[
 S=|a|_1+|b|_1+1,\quad Q=|a|_1|b|_1,\quad
 P=SR+|c|_1,\quad N=Q+P.
\]
The Gauss equation \cite{dlmf-gauss} becomes a first-order system in $Y=(y,y')$:
\[
 \frac{dY}{dz}=A(z)Y,\qquad
 A(z)=\begin{pmatrix}
 0&1\\[1mm]
 \displaystyle\frac{ab}{d(z)}&
 \displaystyle\frac{(a+b+1)z-c}{d(z)}
 \end{pmatrix}.
\]
On the segment the coefficient matrix is $B(t)=hA(z(t))$. Define natural
numbers by exact rational ceilings:
\begin{align}
 K&=\left\lceil T(1+ND_*)\right\rceil,\label{eq:gauss-K}\\
 L&=\left\lceil T^2\bigl(SD_*+ND_*^2(1+2R)\bigr)\right\rceil.
 \label{eq:gauss-L}
\end{align}

\begin{proposition}[Produced Gauss transport data]\label{prop:gauss-transport-data}
The matrices $B(t)$ satisfy \eqref{eq:ode-KL} with
\eqref{eq:gauss-K}--\eqref{eq:gauss-L}. Consequently a certified rational
segment avoiding $0,1$, together with initial values, produces the regular
name and explicit precision schedule of \eqref{eq:ode-demands}.
\end{proposition}
\begin{proof}
The first row of $B$ has row norm $T$. The second has row norm at most
$TND_*$ by \eqref{eq:ode-inverse-bound}; their maximum is bounded by
$T(1+ND_*)$. For differences, the first row is constant. For the second
row use \eqref{eq:ode-inverse-lip}; the two constant numerator bounds sum
to $N$, and the varying numerator has Lipschitz bound $S$ in $z$.
Thus the row difference is at most
$T\bigl(SD_*+ND_*^2(1+2R)\bigr)|z-w|_1$.
Since $|z(t)-z(s)|_1=T|t-s|$, \eqref{eq:gauss-L} follows.
\end{proof}

\begin{proposition}[Producing coefficient-name stability from parameter names]
\label{prop:gauss-parameter-stability}
Suppose two parameter triples $(a,b,c)$ and $(\widetilde a,\widetilde b,
\widetilde c)$ have all six $|\cdot|_1$ norms at most a natural $A$, and
corresponding entries differ by at most $\varepsilon$. For the same segment,
\[
 \lVert B(t)-\widetilde B(t)\rVert\le J\varepsilon,
 \qquad
 J=\left\lceil TD_*(2A+2R+1)\right\rceil.
\]
The bounds \eqref{eq:gauss-K}--\eqref{eq:gauss-L} hold uniformly for every
such triple after using
$S=2A+1$, $Q=A^2$, $P=(2A+1)R+A$, and $N=Q+P$.
\end{proposition}
\begin{proof}
The denominator depends only on the fixed path. The identity
$ab-\widetilde a\widetilde b=a(b-\widetilde b)+
\widetilde b(a-\widetilde a)$ bounds the first numerator difference by
$2A\varepsilon$. The second numerator difference is
$((a-\widetilde a)+(b-\widetilde b))z-(c-\widetilde c)$ and has norm at most
$(2R+1)\varepsilon$. Add the two row contributions and multiply by $TD_*$.
The uniform majorants follow directly from the common parameter bound.
\end{proof}

This proposition supplies an actual coefficient-name producer. If $a_r,b_r,c_r$
are three particular PR regular names with all approximants bounded by $A$,
substitution in the rational Gauss matrix gives PR kernels with the proved
common $K,L$ and the proved comparison bound $J(2^{-r}+2^{-s})$.
For a single segment, a safe joint schedule for the parameter names, the
initial name, and the Euler mesh uses
\[
 c_{\mathrm{par}}=\ell(8G^2U\max(1,J)),\qquad
 c_v=\ell(8G),\qquad c_m=\ell(8\max(1,C)).
\]
At output precision $p$, query all three parameters at $p+c_{\mathrm{par}}$,
the initial value at $p+c_v$, and use $2^{p+c_m}$ cells.
Comparing two outputs uses, in order, the mesh error, the initial-data error,
and the parameter error. Their upper bounds for $q\ge p$ are respectively
$2^{-p}/8$, $(2^{-p}+2^{-q})/8$, and
$(2^{-p}+2^{-q})/8$, whose sum is at most $5\cdot2^{-p}/8$.
Thus parameter stability is not left as an uninhabited analytic premise.

This ordinary-point transport works for resonant parameters too. Producing
the particular initial values of ${}_2F_1(a,b;c;z)$ from its regular series
still requires the series packet's parameter and denominator conditions.
The transport does not establish gamma connection coefficients, contour
bounds, or modular inversion. Those are different producer obligations.

\subsection{A whole finite path with one closed bound}

Let a finite rational polygonal path have $m$ certified edges, with data
$K_j,L_j$ and $G_j=3^{K_j}$. For $m=0$ the output is the input. For $m>0$
put
\begin{equation}\label{eq:ode-path-constants}
 G_{\mathrm{tot}}=\prod_{j<m}G_j,\qquad
 C_{\mathrm{path}}=
 U G_{\mathrm{tot}}\sum_{j<m}G_j(L_j+K_j^2).
\end{equation}
Apply the Euler program with the same $n=2^r$ subdivisions on each edge,
feeding the resulting rational endpoint into the next edge.

\begin{theorem}[Finite-path endpoint name]\label{thm:ode-path-name}
For rational initial data, consecutive global refinements have discrepancy
at most $C_{\mathrm{path}}/(2\cdot2^r)$, and the tail from level $r$ to any
finer level is at most $C_{\mathrm{path}}2^{-r}$. Thus a finite path has a
primitive-recursive endpoint name. For named initial data, valid shifts are
$\ell(4G_{\mathrm{tot}})$ for the input and
$\ell(4\max(1,C_{\mathrm{path}}))$ for the mesh exponent.
\end{theorem}
\begin{proof}
Let $E_{j,n}$ be the linear endpoint map of the $j$th Euler program.
Its operator norm is at most $G_j$. The proof of
Theorem~\ref{thm:ode-produced-refinement}, applied to the matrix initial value $I$ with row-sum norm at most one, gives
\[
 \lVert E_{j,n}-E_{j,2n}\rVert
 \le\frac{G_j^2(L_j+K_j^2)}{2n}.
\]
The difference of two finite products is the sum obtained by changing one
factor at a time. The $j$th term, acting on $v$, has norm at most
\[
 U\left(\prod_{\ell>j}G_\ell\right)
 \frac{G_j^2(L_j+K_j^2)}{2n}
 \left(\prod_{\ell<j}G_\ell\right)
 =\frac{UG_{\mathrm{tot}}G_j(L_j+K_j^2)}{2n}.
\]
Add these finitely many bounds and then telescope dyadic levels. The product
also has initial-data sensitivity $G_{\mathrm{tot}}$, giving the two named
input shifts by the proof of \eqref{eq:ode-demands}.
\end{proof}

For named Gauss parameters shared along the entire path, let $J_j$ be the
produced parameter Lipschitz constant on edge $j$, and put
\[
 S_{\mathrm{path}}=UG_{\mathrm{tot}}\sum_{j<m}G_jJ_j.
\]
The same one-factor-at-a-time product comparison, now using
Lemma~\ref{lem:ode-stability} rather than the mesh estimate, bounds the
endpoint change by $S_{\mathrm{path}}\varepsilon$ when the three parameter
changes are at most $\varepsilon$. A complete joint path schedule therefore
uses shifts $\ell(8\max(1,S_{\mathrm{path}}))$ for the parameter names,
$\ell(8G_{\mathrm{tot}})$ for the initial name, and
$\ell(8\max(1,C_{\mathrm{path}}))$ for the common mesh exponent. All constants
are finite products and sums over the supplied path.

\subsection{The same producer for higher-order equations and Meijer \texorpdfstring{$G$}{G}}

The construction is not restricted to two-dimensional systems. A scalar equation
\[
 a_q(z)y^{(q)}+\sum_{j<q}a_j(z)y^{(j)}=0,\qquad q\ge1,
\]
with polynomial coefficients becomes a rational companion system for
$(y,y',\ldots,y^{(q-1)})$ whenever a finite certificate bounds the reciprocal
of $a_q$ on the path. The superdiagonal entries are one and the last row is
$-a_j/a_q$. Polynomial coefficient sums and the finite power-difference
identity compute bounds for each numerator and its Lipschitz constant.
The reciprocal-difference identity supplies the remaining denominator
estimate. Thus this is one fixed PR kernel with numerical order and
coefficient lists, not a separate implementation for every equation.

For a concrete class, let $\theta=z\,d/dz$ and
$\kappa=(-1)^{p-m-n}\in\{-1,1\}$. The Meijer $G$ equation \cite{dlmf-meijer} is
\[
 \left(\kappa z\prod_{i=1}^{p}(\theta-a_i+1)
       -\prod_{j=1}^{q}(\theta-b_j)\right)y=0
.
\]
Here $p,q,m,n$ are natural numbers with $0\le m\le q$, $0\le n\le p$; take $p\le q$ and $q\ge1$. Expand the two finite products as
polynomials in $\theta$. The identity
\[
 \theta^k=\sum_{j=0}^k S(k,j)z^j\frac{d^j}{dz^j},\qquad
 S(k+1,j)=jS(k,j)+S(k,j-1),
\]
with $S(0,0)=1$ and out-of-range entries zero, converts them to ordinary
derivatives. To prove it, apply $z\,d/dz$ to each summand and collect the
two resulting coefficients. This finite recurrence is PR and its coefficient
identity has a QF-induction proof. The highest coefficient is
\[
 a_q(z)=
 \begin{cases}-z^q,&p<q,\\z^q(\kappa z-1),&p=q.\end{cases}
\]
For a rational segment the exact clearance procedure above therefore needs
only the point $0$, and additionally $\kappa$ when $p=q$.

Here are explicit companion bounds. Let $R$ bound $|z|_1$ on the segment,
$T=|h|_1$, and $\mu_s>0$ be its minimum squared distance to each needed
singularity. A reciprocal bound is
\[
 D_q=(R/\mu_0)^q\quad(p<q),\qquad
 D_q=(R/\mu_0)^q(R+1)/\mu_\kappa\quad(p=q).
\]
Write $a_j(z)=\sum_k c_{jk}z^k$, and compute
$P_j=\sum_k|c_{jk}|_1R^k$ and
$L_j=\sum_{k\ge1}k|c_{jk}|_1R^{k-1}$, including $j=q$.
Then $|a_j(z)|_1\le P_j$, its difference is bounded by
$L_j|z-w|_1$, and the reciprocal difference is bounded by
$D_q^2 L_q|z-w|_1$. Consequently valid natural constants are
\[
 K=\left\lceil T\left(1+D_q\sum_{j<q}P_j\right)\right\rceil,
\]
\[
 L=\left\lceil T^2\left(D_q\sum_{j<q}L_j+
                 D_q^2 L_q\sum_{j<q}P_j\right)\right\rceil.
\]
These formulas, the finite expansion, and the clearance tests produce all
the rational transport data from finite rational parameters. Repeated
parameters do not invalidate an ordinary-point companion equation.
They can, however, change the difficulty of constructing the initial jet
of the particular Meijer $G$ branch. This section supplies continuation
\emph{from a supplied certified jet}; it does not identify that jet with a
Mellin--Barnes contour integral or prove a connection formula.

\subsection{Integral observations, uniqueness, and the semantic boundary}

A Cauchy estimate alone should not be advertised as an ODE theorem. Here the
finite output also satisfies the required integral observations. Let
$F_n(t)=B(t)X_n(t)$. From \eqref{eq:ode-KL} and
Lemma~\ref{lem:ode-polybound},
\[
 \lVert F_n(t)-F_n(s)\rVert\le H|t-s|.
\]
On the $j$th Euler cell the polygon slope is the constant
$Q_n(t)=B(j/n)X_{n,j}$, and
\begin{equation}\label{eq:ode-local-residual}
 \lVert F_n(t)-Q_n(t)\rVert\le H/n.
\end{equation}
The finite integral of $Q_n$ is exactly $X_n(t)-v$, by the defining recurrence
and the last partial cell. For an independent positive mesh number $M$, define
\[
 R_M(F_n;t)=\frac{t}{M}\sum_{i<M}F_n(it/M).
\]
A common rational refinement of the Euler cells and the $M$ quadrature cells
proves
\begin{equation}\label{eq:ode-integral-observation}
 \left\lVert X_n(t)-v-R_M(F_n;t)\right\rVert
 \le H/n+H/M.
\end{equation}
Indeed, replace $F_n$ on each quadrature cell by its left value. Its variation
is at most $H/M$, and the sum of all cell lengths is at most one. Next use
\eqref{eq:ode-local-residual} on the common refinement to replace $F_n$ by
$Q_n$; that contributes at most $H/n$. This is a proof about two finite sums,
not an invocation of integration or completeness inside iPRA. Choosing
$n,M$ to make the two rational upper bounds at most a prescribed dyadic
error produces the integral-equation observation contract.

The uniform $KY$-Lipschitz bound on all Euler polygons passes to their finite observation laws. Thus the pointwise names are spatially coherent and represent one uniformly continuous function. In a semantic interpretation of regular names, the same estimates identify
the output with the solution of
\[
 X(t)=v+\int_0^tB(s)X(s)\,ds.
\]
For completeness, uniqueness has a constructive quantitative proof. Suppose
a second solution $Z$ has the same initial value and a bound
$\lVert Z(t)\rVert\le V$. Its integral equation implies that it is
$KV$-Lipschitz. Hence $B(t)Z(t)$ is $(L+K^2)V$-Lipschitz. On a mesh cell,
subtract its left rectangular contribution from its integral; the remaining
norm is at most $(L+K^2)V/(2n^2)$. Comparison with Euler and finite Gronwall
therefore gives
\[
 \lVert Z(j/n)-X_{n,j}\rVert
 \le\frac{G(L+K^2)V}{2n}.
\]
The interpolation estimate gives convergence off the nodes as well. For
every requested precision choose the displayed finite $n$ and the necessary
finite integral observations; the two solution names must agree at that
precision. No zero decision or infinite subsequence is selected.

Splitting an edge at a rational fraction $\theta$ replaces $B(t)$ on the
first subinterval by $\theta B(\theta t)$ and on the second by
$(1-\theta)B(\theta+(1-\theta)t)$. The corresponding bounds scale as
$K\mapsto\theta K$, $L\mapsto\theta^2L$ before rounding up. Adding the two
finite integral observations recovers the original one. The uniqueness
estimate proves equality of the endpoint names under rational subdivision.
This does not imply equality for different paths around singularities;
monodromy is compatible with every theorem above.

The formal object statement quantifies natural codes and precision variables
for particular compiled kernels. It does not quantify all real functions.
The semantic statement that such names denote classical or constructive
solutions is a separate interpretation theorem. Keeping this boundary visible
both respects iPRA's language and preserves the substantive analytic content.

\subsection{Lean implementation plan at the actual source boundary}

The following dependency order produces concrete endpoints instead of a
new layer of uninhabited contracts.

\begin{enumerate}
\item Encode signed rational arrays, positive denominators, finite maximum
and row-sum norms, and guarded inverse. Establish the displayed finite norm
inequalities as object-theory lemmas.
\item Define the primitive-recursive Euler node and interpolation programs.
Derive their equations, totality of the codes, and denominator positivity.
\item Prove the node growth and polygon Lipschitz bounds. Reuse the existing
\nolinkurl{DiscreteGronwall.growth}; induction invariants are denominator-cleared
quantifier-free inequalities.
\item Prove \eqref{eq:ode-residual-identity} by splitting the fine-node index
into its parity cases. Bound its two terms to obtain
\eqref{eq:ode-produced-recurrence}. This is the central new producer lemma.
\item Instantiate \texttt{OdeRefinement.node\_bound} and
\texttt{geometric\_cauchy}, then prove the affine interpolation estimate and
the explicit $\ell$-based name schedule.
\item Prove initial-data and coefficient stability. Expose the exact affine
precision demands to the already existing precision/cost graph interface.
\item Implement segment clearance by the three rational clamping cases and
bind the Gauss matrix to \eqref{eq:gauss-K}--\eqref{eq:gauss-L}.
\item Prove finite product telescoping and obtain
\eqref{eq:ode-path-constants} uniformly in coded path length. Connect the
initial-value producer supplied by the particular special-function packet.
\item Formalize the common-refinement sum proof of
\eqref{eq:ode-integral-observation}. Treat the semantic solution
interpretation and the object-level finite observation theorem as separately
named declarations.
\end{enumerate}

For a fixed PR coefficient kernel, the decisive eventual theorem has literal
judgment \texttt{Prf PRAAx}: from denominator guards and the rational
coefficient bounds it derives the coded Euler refinement inequality. A
second literal \texttt{Prf PRAAx} theorem assembles the named outputs. A
host Lean theorem about \texttt{Real}, a numerical run, or successful
elaboration of an interface structure does not substitute for either
endpoint.

\section{Lean architecture and a concrete route to object-level proofs}\label{sec:lean}

\subsection{Two layers with different purposes}

The supplied \texttt{lean/} directory is a specification and reference layer. It uses rational numbers, finite vectors, lists and matrices in ordinary Lean. \nolinkurl{FeedbackReference.lean} contains actual candidate proof scripts for finite contraction and tail lemmas; the other blueprint files define the numerical mathematics and name the remaining propositions as \texttt{Prop}-valued goals. None of these files was compiled in the delivery environment. No new axiom or admitted theorem is used to present an unfinished target as a proof.

The second layer must be implemented in ProveIt's \nolinkurl{IPRADerivations}. It must return actual values of \texttt{Prf PRAAx}, using the checked numerical encodings. Its final theorems quantify over the numerical input data and precision, and retain the original certificate guard. An ordinary Lean theorem about rational matrices is a reference for this layer, not an automatic transfer theorem.

The baseline repository specifies Lean \texttt{v4.32.0} and Mathlib commit
\nolinkurl{81a5d257c8e410db227a6665ed08f64fea08e997}. The included Lake configuration pins that Mathlib commit for subsequent checking. The precise baseline and source hashes appear in the provenance manifest. Reproduction should first build the small reference files, then port one object endpoint at a time.

\subsection{Definitions supplied in the reference layer}

\begin{longtable}{>{\raggedright\arraybackslash}p{.30\textwidth}>{\raggedright\arraybackslash}p{.62\textwidth}}
\toprule File & Principal definitions and intended role\\\midrule\endhead
\nolinkurl{FeedbackReference.lean} & Rational finite vectors and matrices; \texttt{rowApply}; componentwise \texttt{Within}; positive half-gain certificates; invariant-domain maps; primitive iteration; finite maximum witnesses. Candidate proofs establish row domination, half contraction, increments, finite tails and residual comparison.\\
\nolinkurl{PolynomialBlueprint.lean} & Finite rational polynomial terms; coefficient algebra and expansion; preconditioned update; coefficient gain matrix; checked finite certificate; iteration, residuals and explicit precision goals. The theorem obligations follow Sections~\ref{sec:polynomial} and \ref{sec:factors}.\\
\nolinkurl{FactorBlueprint.lean} & Finite coefficient convolution, Sylvester map and inverse checker, norm bound, actual factor iteration, residual and precision clock. Its unproved goals express the complete finite lifting theorem.\\
\nolinkurl{ClusterBlueprint.lean} & Numerical finite graph and matching predicates; neighborhood counts and defects; bounded public matching selector; weighted-cover and stable-locator contracts. Unproved goals correspond to the constructive proof of Section~\ref{sec:cluster-transport}.\\
\nolinkurl{OdeBlueprint.lean} & Rational finite norms, coefficient contracts, Euler nodes, interpolation, coarse residuals, odd-node residual expression, bounded precision clock, named outputs and quadrature sums. Unproved goals isolate the produced recurrence, refinement, stability, clearance and observation adapters.\\
\bottomrule
\end{longtable}

The common feedback definitions have the following shape. The complete code, rather than this schematic extract, is in the archive.
\begin{lstlisting}
abbrev Vec (m : Nat) := Fin m -> Rat
abbrev Mat (m : Nat) := Fin m -> Fin m -> Rat

-- A finite numerical assertion: no internal real-number sort.
def Within (w : Vec m) (r : Rat) (x y : Vec m) : Prop :=
  forall i, abs (x i - y i) <= r * w i

structure HalfCertificate (m : Nat) where
  A : Mat m
  w : Vec m
  weight_pos : forall i, 0 < w i
  entry_nonneg : forall i j, 0 <= A i j
  row_bound : forall i, 2 * rowApply A w i <= w i
\end{lstlisting}
Here a Lean structure packages hypotheses. It does not mean a mere file containing numbers certifies itself. The numerical implementation has a separate PR checker that verifies exactly these inequalities and shape conditions, and a proof that acceptance supplies the corresponding object-level formulas.

\subsection{The actual numerical programs to internalize}

Use the existing signed-rational representation, sequence encoder, length and entry operations. For each task, define one fixed PR program with the input code packed into its parameter. Keep numerical dimension and precision internal.

\begin{enumerate}[leftmargin=*,itemsep=3pt]
\item \textbf{Feedback.} Encode matrix, weights, initial state and optional rounding schedule. The selected update is one fixed PR template. Store the orbit by primitive recursion; a finite matrix kernel is evaluated by bounded sums.
\item \textbf{Polynomial systems.} Encode the monomial arrays and all exponents, generate the translated update polynomial, combine equal exponent lists by bounded searches, and compute \(W\) by \eqref{eq:polygain}. The checker verifies \eqref{eq:polycert} and the inverse equation. Runtime polynomial syntax is allowed because its evaluator is explicitly PR.
\item \textbf{Factor lifting.} Construct the Sylvester map by indexed shifts, check a rational inverse, compute its column-sum bound, verify \eqref{eq:factorcertificate}, and iterate \eqref{eq:factoriteration}. Inverse production can use a fixed bounded elimination routine. Its correctness is a separate finite matrix theorem; a user-supplied inverse still must be checked.
\item \textbf{Matching.} Use the Boolean graph, bounded queue, marks, parent array and partial permutation from the constructive proof. Alternatively, after proving existence, a bounded scan through all length-\(N\) base-\((N+1)\) lists selects the first matching; the explicit scan bound is \((N+1)^N\).
\item \textbf{ODEs.} Encode the Euler node list and rational interpolation program. For Gauss or Meijer parameters, the rational kernel is the concrete companion formula. Matrix size, polynomial order, number of segments, mesh and input precision are natural variables.
\end{enumerate}

Every invalid input receives a fixed totalized output. Acceptance guards stay in theorem statements. The resulting programs do not use unrestricted minimization, a type containing all PR functions, or an external answer whose correctness is checked only after termination.

\subsection{The QF invariant for each main proof}

For feedback let \(\operatorname{Orbit}(c,n)\) be the coded orbit and let
\(\operatorname{Inc}(c,n)\) be the comparison bit for
\[
 2^n\abs{x_{n+1,i}-x_{n,i}}\le B w_i\quad\text{for every }i<m(c).
\]
Its comparison is performed after clearing the stored positive denominators. The induction formula is a conjunction of \(\operatorname{DomainBit}(c,n)=1\) and \(\operatorname{Inc}(c,n)=1\), not a statement quantifying over all future times.

For finite tails, fix the initial index \(n\) as a parameter and induct on the finite offset \(j\). The comparison bit expresses the sharp geometric bound in \eqref{eq:feedback-tail}. Multiplication by \(2^{n+j}\) removes dyadic division. For variable dimension the universal coordinate condition is one bounded product of bits.

For polynomial compilation, the invariant is coefficient equality of the partial expansion and its evaluation at an arbitrary numerical point. An outer fold over monomials and an inner fold over coordinates prove the translation identity. The monomial difference proof is a finite telescoping identity plus a bounded sum of absolute-value inequalities. One must prove this for the actual list decoder, including zero coefficients and missing entries, rather than relying on a host-level function equality.

For factor lifting, retain three numerical facts in the invariant: lengths and monicity; the coefficient ball bound; and the scaled successive-displacement bound. The residual identity is a finite polynomial identity derived at an arbitrary iteration index. No induction over the assertion that a limiting factorization exists is required.

For matching, Booleanize the partial-bijection equations, discovered-parent conditions and closure assertions. The number of marked vertices supplies the finite search bound. A failed augmentation outputs its actual neighbor mask and cardinality comparison. The two alternatives are generated by decidable finite tests; no classical Hall theorem is imported into the object theory.

For Euler refinement, pair the node-bound invariant with the finite comparison recurrence. Parity splits the residual into its exact even and odd formulas. Prove the interpolation identity using quotient/remainder at a rational argument. The comparison kernel then instantiates the existing \texttt{OdeRefinement.node\_bound}; the tail instantiates its geometric summation. For runtime vector dimension, all coordinate and norm comparisons are reflected into finite Boolean folds.

\subsection{Representative closed object sentences}

The following notation denotes explicit formulas assembled from the actual PR symbols. It is not an additional schema of axioms. Let \(\mathrm{Inv}\), \(\mathrm{Close}\) and \(\mathrm{Res}\) be the reflected bits for the polynomial invariant, tail and residual in Theorem~\ref{thm:polysolver}. The desired endpoint is
\begin{equation}\label{eq:objectpoly}
 \ipra\vdash\forall c,n,j\;
 \bigl(\mathrm{Cert}(c)=1\to
   \mathrm{Inv}(c,n)=1\wedge\mathrm{Close}(c,n,j)=1
       \wedge\mathrm{Res}(c,n)=1\bigr).
\end{equation}
The matrix sizes, polynomial degrees and all input arrays are decoded from \(c\); \(n,j\) are internal naturals. The certificate formula contains no assertion about an already constructed solution.

For matching the endpoint is the matching-or-defect disjunction stated in Section~\ref{sec:cluster-transport}. Its coverage corollary retains the finite occurrence/count assumptions and existentially seals any private incidence witnesses. For rational ODE kernels the endpoint has the form
\[
 \ipra\vdash\forall c,n,t\;
 \bigl(\mathrm{AdmissibleKernel}(c)\wedge n>0\wedge0\le t\le1
 \to 2n\,\norm{X_n(t)-X_{2n}(t)}\le C(c)\bigr),
\]
where the norm inequality is the actual denominator-cleared rational comparison. For generic fixed kernels, the coefficient contracts are explicit antecedents; for the concrete companion kernels, those antecedents are proved from the clearance and finite parameter data.

At the Lean level the endpoint type must be an actual sequent, schematically
\begin{lstlisting}
([] : List (Fm PTm 0)) |- [PRAAx] univ 3 polynomialEndpoint
\end{lstlisting}
with the repository's genuine notation and scoping machinery. The archive does not declare this theorem with an empty proof. It records the exact formula obligations and the implementation order needed to obtain it.

\subsection{Reuse and acceptance criteria}

\begin{longtable}{>{\raggedright\arraybackslash}p{.30\textwidth}>{\raggedright\arraybackslash}p{.62\textwidth}}
\toprule Existing modules & Reuse in the proposed port\\\midrule\endhead
\nolinkurl{Core/Pairing}, \nolinkurl{Foundations/SequenceComposition}, \nolinkurl{MasterCertificate} & Length-checked arrays, append and entry identities; parameters and runtime dimensions.\\
\texttt{Core/Order}, \nolinkurl{BoundedSyntax}, \nolinkurl{BoundedReflection} & Rational comparison bits, bounded quantifier reflection, QF invariants.\\
\nolinkurl{Analysis/RationalArithmetic}, rational sum modules & Exact rational addition, multiplication, finite sums and positive-denominator reasoning.\\
\nolinkurl{Core/SumLemmas}, \nolinkurl{Core/PolynomialInequalities} & Finite norm inequalities, row domination, polynomial and convolution bounds.\\
\nolinkurl{DiscreteGronwall}, \nolinkurl{OdeRefinement} & Already checked integer growth, recurrence estimates and geometric Cauchy arithmetic.\\
\nolinkurl{Algebra/HenselStep} & Ring identities for the correction and residual, after connecting their finite coefficient instances.\\
\nolinkurl{PrecisionCostDAGSchedule}, \nolinkurl{OpenContracts} & Final scheduling and composition after the new producer hypotheses are discharged.\\
\bottomrule
\end{longtable}

For each new object module, build it at the pinned provider versions, inspect the exact theorem type, and run \texttt{\#print axioms}. Absence of \texttt{sorryAx} is necessary, but the more important condition is that the target is an actual \texttt{Prf PRAAx} value with the intended original hypotheses. A theorem about an evaluator, an assumed local estimate, or an abstract proof-bearing structure must not be reported as the completed application endpoint.

The recommended first port is the scalar/weighted finite comparison in \nolinkurl{FeedbackObjectPort.md}, followed by the actual polynomial difference compiler. The factor example supplies a small, exact end-to-end milestone. The ODE port can proceed independently using the existing Gronwall backend; the matching port has a separate finite combinatorial invariant. This division exposes the remaining work without circular module dependencies.

\section{Resolved interfaces, limitations, and further mathematical uses}\label{sec:boundaries}

\subsection{What has been established}

The framework now has concrete producers in two broad classes. A finite polynomial certificate produces a local implicit solution with a residual and stability modulus. A rational linear differential system with coefficient bounds produces a coherent solution observation along a supplied finite path; Gauss and Meijer companion systems have explicit coefficient and clearance producers. These are sufficiently general to support many algebraic and special-function applications without proving a new existence theorem every time.

The root-cluster results resolve the specific mathematical ambiguities in FIBER G1 and G2. With complete occurrence covers, a bounded matching algorithm obtains the sharp sum-of-radii matching bound. With only component balance, the weaker component-diameter bound is the correct theorem. Stable locators or stable factor coefficients replace unguarded lexicographic root labels. Producing the multiplicity covers from arbitrary polynomial coefficients remains a separate task.

The new human proofs deliberately end in finite statements strong enough to be useful: a computable observation, its admissibility guard, a finite error law, and an explicit precision demand. The new object-language Lean proofs remain implementation work. This is a substantive mathematical extension and a formalization blueprint, not a report of newly checked endpoints.

\subsection{A further correction: spanning is not a basis certificate}

The same discipline clarifies the FIBER G3 finite-algebra boundary. A list spanning a finite algebra does not make the trace of an arbitrary lifted multiplication matrix intrinsic.

For a concrete counterexample, let \(A=\Q[X]/(X)\), and use the redundant spanning list \((1,1)\). Multiplication by the residue class of \(X\) is zero. For any rationals \(a,b\), the matrix
\[
 M=\begin{pmatrix}a&b\\-a&-b\end{pmatrix}
\]
represents that zero action on the spanning list, because each column sums to zero. Yet \(\operatorname{tr}M=a-b\) is arbitrary. The true one-dimensional algebra has trace zero for this endomorphism. Thus neither an intrinsic trace nor a characteristic polynomial of the quotient follows from spanning alone.

There is a useful positive replacement. Let \(b_1,\ldots,b_N\) span a commutative algebra over \(\Q\), let \(x\) be an algebra element, and suppose a matrix \(M\) satisfies the finite relations
\(xb_j=\sum_iM_{ij}b_i\). If the finite Cayley--Hamilton identity
\(\chi_M(M)=0\) is available, then \(\chi_M(x)b_j=0\) for every \(j\): first prove the assertion for each power by induction, then take the finite linear combination defining \(\chi_M\). Since the list spans the unit, there are coefficients \(c_j\) with \(1=\sum_jc_jb_j\), and
\[
 \chi_M(x)=\sum_j c_j\chi_M(x)b_j=0.
\]
Thus the matrix supplies a valid annihilating polynomial. It need not supply intrinsic spectral multiplicities or trace data. To obtain those, require an explicit basis/freeness certificate, or a proved invariant under the particular change of presentation. In iPRA all relations and coefficients in this argument must be supplied by finite arithmetic certificates; a general algebra object is only an external explanatory shorthand.

\subsection{Why these methods do not settle arbitrary difficult theorems}

There are several distinct possible obstructions, and conflating them would hide the real work.

\textbf{An absent quantitative producer.} A proof may assert that a favorable radius, separation, height bound or finite cover exists, without constructing a PR bound or a verified finite certificate. The weighted contraction theorem does not create such data. The polynomial and ODE applications are useful precisely because they compute the needed majorants and acceptance tests for their declared input class.

\textbf{A discontinuous requested observation.} Exact sign of an unrestricted real name, arbitrary complex-root ordering, and choice of a unique branch through a singularity can fail finite-use continuity. This report changes neither the input representation nor the theorem silently: it states a gap, locator or branch certificate when one is needed. The repository already contains a finite-use sign obstruction \cite{repo-sign}.

\textbf{A nonuniform termination argument.} Showing that an algorithm eventually succeeds does not itself exhibit a primitive recursive clock. The clocks here are explicit: geometric precision shifts, bounded graph exploration, finite Gaussian elimination, and Euler mesh sizes. If a future application repeatedly subdivides until a certificate accepts, it must supply a prior bound or a separate termination theorem in iPRA. Calling the loop ``adaptive'' does not discharge that obligation.

\textbf{A missing identity with the intended mathematical object.} A solution of a companion ODE with some initial jet is not automatically the particular special function named by a contour integral. Gamma connection formulas, residue expansions and initial jets keep their own proofs. Likewise, a small polynomial residual is not a complete root-multiplicity certificate. The present framework exposes these exact interfaces and supplies some of their producers; it does not rename an unproved hypothesis as a theorem.

\textbf{A stronger foundational requirement.} These techniques do not prove general Siegel or Faltings theorems in iPRA, nor do they prove such statements impossible there. A missing effective bound is evidence about the current argument, not by itself an independence theorem. One must state the precise arithmetical version, give its actual finite witnesses or negative interpretation, and analyze its induction/choice requirements separately.

\subsection{How to reuse the framework in another packet}

For a new application, first specify the finite observable. Choose a representation whose equality and approximation claims have a finite arithmetic meaning. Then write the actual evaluator or producer and state its domain. Identify the numerical quantities that control perturbations: a coefficient majorant matrix, an inverse bound, a path clearance, a coverage count, or a locator gap.

Next derive the appropriate local inequality from the input data. If the local errors depend cyclically on one another, certify a positive weighted gain instead of composing the same contract around a cycle. Compute the resource/precision schedule from that certificate. Prove the finite tail, residual, or matching statement with a named QF invariant. Finally show that the resulting observation is independent of representation choices and agrees with the particular mathematical object in the packet.

This sequence separates reusable arithmetic from the substantive mathematics of each application. It also makes progress measurable: a formerly assumed local contract is replaced by a concrete producer and a finite proof, while any remaining source obligation stays visible in the endpoint.

\appendix
\section{Artifacts, checks, and reproducibility}\label{sec:artifacts}

The archive contains the complete LaTeX source and this compiled PDF, the Lean reference files and implementation plan, four exact-rational reference programs, a batch verification driver, a source manifest, and a claim ledger. The source is designed to build with an ordinary LaTeX installation; the Python reference programs require only the standard library.

The four computational checks have deliberately limited but useful roles.
\begin{itemize}[leftmargin=*,itemsep=3pt]
\item The polynomial check compares the generated shifted polynomial against direct evaluation, tests the coefficient-derived gain inequality, verifies the coupled nonlinear orbit estimates, checks the exact Sylvester inverse and residual identity, and compares factor iterates under changed coefficients.
\item The feedback check exercises a cyclic positive matrix whose unweighted row bound is too large, a nonlinear invariant-box example, finite residual comparison and perturbation, the zero-dimensional case and a rejected gain certificate.
\item The cluster check exhausts every bipartite Boolean graph with equal side size at most three, validating either the returned permutation or the returned Hall defect. It also verifies the exact component counterexample and the locator/lexicographic examples.
\item The ODE check verifies exact even/odd residual identities, refinement inequalities at rational samples, coefficient perturbation, exact segment clearance, and the Gauss coefficient and parameter majorants on complex-rational paths.
\end{itemize}

All reported inequalities in these programs use exact rational arithmetic. Decimal values, when present, are explicitly for orientation. Sampling is never used in the mathematical proof of a universal bound. For example, the ODE Lipschitz theorem follows from the finite algebraic estimates in the paper, not from its sampled test points.

From the archive root run
\begin{lstlisting}
python3 code/run_all.py
python3 code/build_report.py
\end{lstlisting}
The first command regenerates the numerical evidence and a summary. The second compiles the LaTeX source and places the resulting report at the archive root, using the local LaTeX engine. It does not download anything. The delivery environment needed local TeX search/format preparation; the build helper supports that reproducible fallback when the installed distribution has packages but no generated format.

To begin checking the Lean reference layer in an environment with Lean installed, run
\begin{lstlisting}
cd lean
lake update
lake build
lake env lean AxiomAudit.lean
\end{lstlisting}
The pinned dependency is recorded in \nolinkurl{lakefile.toml}. These commands are instructions for subsequent checking; they were not executed successfully in the delivery environment, which had no Lean compiler. Any elaboration errors must be repaired before the candidate proofs acquire checked status. Even a successful reference build would leave the separate object-theory port in Section~\ref{sec:lean}.

The claim ledger records each result's assumptions, proof location, executable evidence and Lean status. The source manifest pins the repository commit, inspected paths and byte hashes. This avoids using an evolving branch as if it were a fixed theorem provider, and avoids confusing historical packet status with current derivation status.

\begin{thebibliography}{99}

\small

\bibitem{repo-readme} Vladimir Reshetnikov and ProveIt contributors. \emph{ProveIt: iPRA directory and verification-status guide}. Pinned commit \nolinkurl{3ee850b4c2d6}.

\url{https://github.com/VladimirReshetnikov/ProveIt/blob/3ee850b4c2d650cf01263823fd5415963747c5d0/Logic/PrimitiveRecursiveArithmetic/README.md}.

\bibitem{repo-concise} Vladimir Reshetnikov and ProveIt contributors. \emph{ProveIt: ConciseIPRA systems and theorem names}. Pinned commit \nolinkurl{3ee850b4c2d6}.

\url{https://github.com/VladimirReshetnikov/ProveIt/blob/3ee850b4c2d650cf01263823fd5415963747c5d0/Logic/PrimitiveRecursiveArithmetic/ConciseIPRA/README.md}.

\bibitem{repo-dag} Vladimir Reshetnikov and ProveIt contributors. \emph{ProveIt: actual numerical rational precision-cost DAG and safe schedule}. Pinned commit \nolinkurl{3ee850b4c2d6}.

\url{https://github.com/VladimirReshetnikov/ProveIt/blob/3ee850b4c2d650cf01263823fd5415963747c5d0/Logic/PrimitiveRecursiveArithmetic/Derivations/Lean/IPRADerivations/Foundations/PrecisionCostDAGSchedule.lean}.

\bibitem{repo-fiber} Vladimir Reshetnikov and ProveIt contributors. \emph{ProveIt: FIBER claim audit, including G1, G2 and G3}. Pinned commit \nolinkurl{3ee850b4c2d6}.

\url{https://github.com/VladimirReshetnikov/ProveIt/blob/3ee850b4c2d650cf01263823fd5415963747c5d0/Logic/PrimitiveRecursiveArithmetic/Derivations/FIBER-CLAIMS.md}.

\bibitem{repo-hensel} Vladimir Reshetnikov and ProveIt contributors. \emph{ProveIt: finite ring identities for Hensel steps}. Pinned commit \nolinkurl{3ee850b4c2d6}.

\url{https://github.com/VladimirReshetnikov/ProveIt/blob/3ee850b4c2d650cf01263823fd5415963747c5d0/Logic/PrimitiveRecursiveArithmetic/Derivations/Lean/IPRADerivations/Algebra/HenselStep.lean}.

\bibitem{repo-ode} Vladimir Reshetnikov and ProveIt contributors. \emph{ProveIt: Euler refinement arithmetic and its assumed recurrence}. Pinned commit \nolinkurl{3ee850b4c2d6}.

\url{https://github.com/VladimirReshetnikov/ProveIt/blob/3ee850b4c2d650cf01263823fd5415963747c5d0/Logic/PrimitiveRecursiveArithmetic/Derivations/Lean/IPRADerivations/Foundations/OdeRefinement.lean}.

\bibitem{repo-gronwall} Vladimir Reshetnikov and ProveIt contributors. \emph{ProveIt: discrete Gronwall estimates with rational growth bound}. Pinned commit \nolinkurl{3ee850b4c2d6}.

\url{https://github.com/VladimirReshetnikov/ProveIt/blob/3ee850b4c2d650cf01263823fd5415963747c5d0/Logic/PrimitiveRecursiveArithmetic/Derivations/Lean/IPRADerivations/Foundations/DiscreteGronwall.lean}.

\bibitem{repo-evaluator} Vladimir Reshetnikov and ProveIt contributors. \emph{ProveIt: obstruction to a total primitive recursive universal evaluator}. Pinned commit \nolinkurl{3ee850b4c2d6}.

\url{https://github.com/VladimirReshetnikov/ProveIt/blob/3ee850b4c2d650cf01263823fd5415963747c5d0/Logic/PrimitiveRecursiveArithmetic/Derivations/Lean/IPRADerivations/Foundations/UniversalEvaluatorObstruction.lean}.

\bibitem{repo-sign} Vladimir Reshetnikov and ProveIt contributors. \emph{ProveIt: finite-use obstruction to an unrestricted exact sign selector}. Pinned commit \nolinkurl{3ee850b4c2d6}.

\url{https://github.com/VladimirReshetnikov/ProveIt/blob/3ee850b4c2d650cf01263823fd5415963747c5d0/Logic/PrimitiveRecursiveArithmetic/Derivations/Lean/IPRADerivations/Foundations/FiniteUseSignObstruction.lean}.

\bibitem{dlmf-gauss} NIST Digital Library of Mathematical Functions. \emph{Section 15.10: Hypergeometric Differential Equation}, equation 15.10.1. Consulted 19 November 2026. \url{https://dlmf.nist.gov/15.10}.

\bibitem{dlmf-meijer} NIST Digital Library of Mathematical Functions. \emph{Section 16.21: Differential Equation for the Meijer G-Function}, equation 16.21.1. Consulted 19 November 2026. \url{https://dlmf.nist.gov/16.21}.

\end{thebibliography}

\end{document}
